% Generated mechanically from frozen input inputs/p09/ja/src/spaces-more-groupoids.tex.
% Do not edit this generated file; edit the bound locale source instead.

% OK, start here.
%

\setcounter{chapter}{78}
\chapter{代数空間における亜群の続論}



\phantomsection
\label{spaces-more-groupoids-section-phantom}


\section{はじめに}
\label{spaces-more-groupoids-section-introduction}

\noindent
本章は、代数空間における亜群に関する発展的な話題に充てられる。
結果は代数空間における亜群を用いて述べられるが、読者は $2$-カルテシアン図式
\begin{equation}
\label{spaces-more-groupoids-equation-quotient-stack}
\vcenter{
\xymatrix{
R \ar[r] \ar[d] & U \ar[d] \\
U \ar[r] & [U/R]
}
}
\end{equation}
を念頭に置くべきである。ここで $[U/R]$ は商スタックである。Groupoids in Spaces,
Remark \ref{spaces-groupoids-remark-fundamental-square} を参照せよ。
多くの結果は、この図式を考えることによって動機づけられる。
例えば、Keel と Mori による美しい論文 \cite{K-M} を参照せよ。





\section{記法}
\label{spaces-more-groupoids-section-notation}

\noindent
引き続き Groupoids in Spaces, Section
\ref{spaces-groupoids-section-notation} で導入した規約と記法に従う。





\section{有用な図式}
\label{spaces-more-groupoids-section-diagrams}

\noindent
本章で参照しやすいように、Groupoids in Spaces, Lemmas
\ref{spaces-groupoids-lemma-diagram} および
\ref{spaces-groupoids-lemma-diagram-pull} の結果を簡潔に再掲する。
$S$ をスキームとする。$B$ を $S$ 上の代数空間とする。
$(U, R, s, t, c)$ を $B$ 上の代数空間における亜群とする。可換図式
\begin{equation}
\label{spaces-more-groupoids-equation-diagram}
\vcenter{
\xymatrix{
& U & \\
R \ar[d]_s \ar[ru]^t &
R \times_{s, U, t} R
\ar[l]^-{\text{pr}_0} \ar[d]^{\text{pr}_1} \ar[r]_-c &
R \ar[d]^s \ar[lu]_t \\
U & R \ar[l]_t \ar[r]^s & U
}
}
\end{equation}
において、下側の二つの正方形はファイバー積の正方形である。
さらに、上側の三角形（実際には正方形である）もカルテシアンである。

\medskip\noindent
図式
\begin{equation}
\label{spaces-more-groupoids-equation-pull}
\vcenter{
\xymatrix{
R \times_{t, U, t} R
\ar@<1ex>[r]^-{\text{pr}_1} \ar@<-1ex>[r]_-{\text{pr}_0}
\ar[d]_{\text{pr}_0 \times c \circ (i, 1)} &
R \ar[r]^t \ar[d]^{\text{id}_R} &
U \ar[d]^{\text{id}_U} \\
R \times_{s, U, t} R
\ar@<1ex>[r]^-c \ar@<-1ex>[r]_-{\text{pr}_0} \ar[d]_{\text{pr}_1} &
R \ar[r]^t \ar[d]^s &
U \\
R \ar@<1ex>[r]^s \ar@<-1ex>[r]_t &
U
}
}
\end{equation}
は可換である。上側の二つの行は、与えられた垂直射によって同型である。
左下の二つの正方形はカルテシアンである。







\section{局所構造}
\label{spaces-more-groupoids-section-local}

\noindent
$S$ をスキームとする。
$(U, R, s, t, c, e, i)$ を $S$ 上の代数空間における亜群とする。
$\overline{u}$ を $U$ の幾何学的点とする。この節では、局所環
（Properties of Spaces, Definition
\ref{spaces-properties-definition-etale-local-rings}）
$$
A = \mathcal{O}_{U, \overline{u}}
\quad\text{and}\quad
B = \mathcal{O}_{R, e(\overline{u})}
$$
にどのような構造が得られるかを説明する。ここで用いる規約は、射
$s, t, c, e, i$ が誘導する局所環準同型を、それぞれ対応する文字で表すというものである。
特に、局所環の可換図式
$$
\xymatrix{
A \ar[rd]_t \ar[rrd]^1 \\
& B \ar[r]^e & A \\
A \ar[ru]^s \ar[rru]_1
}
$$
がある。したがって、$I \subset B$ が $e : B \to A$ の核を表すならば、
$B = s(A) \oplus I = t(A) \oplus I$ である。次のように記す。
$$
C = \mathcal{O}_{R \times_{s, U, t} R, (e, e)(\overline{u})}
$$
このとき
$$
C =
(B \otimes_{s, A, t} B)_{\mathfrak m_B \otimes B + B \otimes \mathfrak m_B}^h
$$
である。実際、局所化
$(B \otimes_{s, A, t} B)_{\mathfrak m_B \otimes B + B \otimes \mathfrak m_B}$
の剰余体は分離閉である。$J \subset C$ を $C$ のイデアルで、
$I \otimes B + B \otimes I$ が生成するものとする。
このとき $J$ は局所環準同型
$$
(e, e) : C \longrightarrow A
$$
の核でもある。合成則 $c : R \times_{s, U, t} R \to R$ は環準同型
$$
c : B \longrightarrow C
$$
に対応し、これは $I$ を $J$ に写す。

\begin{lemma}
\label{spaces-more-groupoids-lemma-first-order-structure-c}
写像 $I/I^2 \to J/J^2$ で $c$ が誘導するものは合成
$$
I/I^2 \xrightarrow{(1, 1)} I/I^2 \oplus I/I^2 \to J/J^2
$$
である。ここで第二の射は等式 $J = (I \otimes B + B \otimes I)C$ から得られる。
写像 $i : B \to B$ は写像 $-1 : I/I^2 \to I/I^2$ を誘導する。
\end{lemma}

\begin{proof}
$C$ から別のヘンゼル局所環への局所準同型を記述するには、例えば Algebra, Lemma
\ref{algebra-lemma-henselian-functorial} により、$b_1 \otimes b_2$ の形の元がどう写るかを
述べれば十分である。これを念頭に置くと、二つの標準写像
$$
e_2 : C \to B,\ b_1 \otimes b_2 \mapsto b_1s(e(b_2)),\quad
e_1 : C \to B,\ b_1 \otimes b_2 \mapsto t(e(b_1))b_2
$$
が得られる。これらは埋め込み
$R \to R \times_{s, U, t} R$ で、$r \mapsto (r, e(s(r)))$ および
$r \mapsto (e(t(r)), r)$ によって与えられるものに対応する。
これらの写像が誘導する写像 $J/J^2 \to I/I^2$ は、合わせると補題の写像
$I/I^2 \oplus I/I^2 \to J/J^2$ の逆写像を与える。
% P09-ERR-1346
したがって主張を証明するには、$e_1 \circ c : B \to B$ および
$e_2 \circ c : B \to B$ が恒等写像であることを示せばよい。
これは二つの合成 $R \to R \times_{s, U, t} R \to R$ がいずれも恒等写像であることから従う。

\medskip\noindent
$i$ に関する主張は、$c$ に関する主張と等式 $c \circ (1, i) = e \circ t$ から従う。
細部の一部は省略する。
\end{proof}











\section{切断の亜群}
\label{spaces-more-groupoids-section-groupoid-sections}

\noindent
亜群 $(\text{Ob}, \text{Arrows}, s, t, c, e, i)$ があるとする。
このとき、モノイド $\Gamma$ を構成でき、その元は写像
$\delta : \text{Ob} \to \text{Arrows}$ で $s \circ \delta = \text{id}_{\text{Ob}}$ を満たすものである。
合成は
$$
\delta_1 \circ \delta_2 = c(\delta_1 \circ t \circ \delta_2, \delta_2)
$$
で与えられる。言い換えると、$\Gamma$ の元とは各対象から出る矢を指定する規則 $\delta$ であり、
合成は自然なものである。例えば、明らかな記法のもとで
$$
\vcenter{
\xymatrix{
& \bullet \ar[dl] \\
\bullet \ar@(ul, dl)[] & \bullet \ar[d] \\
& \bullet \ar[lu]
}
}
\quad\circ\quad
\vcenter{
\xymatrix{
& \bullet \ar[d] \\
\bullet \ar[ru] & \bullet \ar[d] \\
& \bullet \ar[lu]
}
}
\quad = \quad\quad
\vcenter{
\xymatrix{
& \bullet \ar@/^/[dd] \\
\bullet \ar@(ul, dl)[] & \bullet \ar[l] \\
& \bullet \ar[lu]
}
}
$$
となる。

\medskip\noindent
同じ手続きは、代数空間における亜群 $(U, R, s, t, c, e, i)$ で
スキーム $S$ 上にあるものに適用できる。すなわち、$\Gamma$ の元として集合
$$
\Gamma = \{\delta : U \to R \mid s \circ \delta = \text{id}_U\}
$$
を取り、合成 $\circ : \Gamma \times \Gamma \to \Gamma$ を上の規則
\begin{equation}
\label{spaces-more-groupoids-equation-composition}
\delta_1 \circ \delta_2 = c(\delta_1 \circ t \circ \delta_2, \delta_2)
\end{equation}
によって与える。単位元は $e \in \Gamma$ で与えられる。
% P09-ERR-1347
亜群 $\Gamma$ は一般には群ではない。実際、逆元をもたない元 $\delta \in \Gamma$ が存在しうる。
$\delta \in \Gamma$ が逆元をもつための必要十分条件が $t \circ \delta$ が $U$ の自己同型であることは明らかであり、
この場合
$\delta^{-1} = i \circ \delta \circ (t \circ \delta)^{-1}$
である。

\medskip\noindent
後で用いるため、部分亜群 $\Gamma_0$ で、$\Gamma$ の切断のうち単位元 $e$ に無限小的に近いものからなるものを考察する。
% P09-ERR-1348
より正確には、$R$-不変閉部分空間 $U_0 \subset U$ が与えられ、$U$ が $U_0$ の一次の厚化であるとする。
$R_0 = s^{-1}(U_0) = t^{-1}(U_0)$ と記し、
$(U_0, R_0, s_0, t_0, c_0, e_0, i_0)$ を対応する代数空間における亜群とする。次のようにおく。
$$
\Gamma_0 = \{\delta \in \Gamma \mid \delta|_{U_0} = e_0\}
$$
$s$ と $t$ が平坦ならば、$\Gamma_0$ のすべての元は可逆である。
実際、$t \circ \delta$ は射 $U \to U$ であって、$\mathcal{O}_{U_0}$ および
$\mathcal{C}_{U_0/U}$ 上で恒等写像を誘導する（Lemma
\ref{spaces-more-groupoids-lemma-idenity-on-conormal}）。したがって、短完全列
$0 \to \mathcal{C}_{U_0/U} \to \mathcal{O}_U \to \mathcal{O}_{U_0} \to 0$
から結論が従う。

\begin{lemma}
\label{spaces-more-groupoids-lemma-idenity-on-conormal}
この節で考えている状況において、$\delta \in \Gamma_0$ および
$f = t \circ \delta : U \to U$ とする。$s, t$ が平坦ならば、標準写像
$\mathcal{C}_{U_0/U} \to \mathcal{C}_{U_0/U}$ で $f$ が誘導するもの
（More on Morphisms of Spaces, Lemma
\ref{spaces-more-morphisms-lemma-conormal-functorial}）は恒等写像である。
\end{lemma}

\begin{proof}
これを見るため、図式 (\ref{spaces-more-groupoids-equation-pull}) の下側を次のように拡張する。
$$
\xymatrix{
Y \ar[r] \ar[d] &
R \times_{s, U, t} R
\ar@<1ex>[r]^-c \ar@<-1ex>[r]_-{\text{pr}_0} \ar[d]_{\text{pr}_1} &
R \ar[r]^t \ar[d]^s &
U \\
U \ar[r]_\delta &
R \ar@<1ex>[r]^s \ar@<-1ex>[r]_t &
U
}
$$
ここで左側の正方形はカルテシアンであり、これを $Y$ の定義とする。$Y$ についてこれ以上知る必要はない。
至る所で $U_0$ へ基底変換することにより、同様の性質をもつ同様の図式が得られる。
示そうとしているのは、$\text{id}_U = s \circ \delta$ および
$f = t \circ \delta$ が余法層上で同じ写像を誘導することである。
$s$ は平坦かつ全射なので、上側の行に沿う二つの合成 $a, b : Y \to R$ について同じことを証明すれば十分である。
$a_0 = b_0$ であり、さらに $a$ と $b$ の一方が同型であることに注意せよ。
実際、$s \circ \delta$ が同型であることを知っている。
したがって、二つの射 $a, b : Y \to R$ は $U$ 上平坦な代数空間の間の射である
（$t : R \to U$ および $t \circ a = t \circ b : Y \to U$ を構造射とする）。
これは所望の主張を含意する。実際、More on Morphisms of Spaces, Lemma
\ref{spaces-more-morphisms-lemma-conormal-functorial-more} の合成との両立性により、二つの写像
$a_0^*\mathcal{C}_{R_0/R} \to \mathcal{C}_{Y_0/Y}$ はいずれも可換図式
$$
\xymatrix{
a_0^*\mathcal{C}_{R_0/R} \ar[rr] & & \mathcal{C}_{Y_0/Y} \\
a_0^*t_0^*\mathcal{C}_{U_0/U} \ar[u] \ar@{=}[rr] & &
(t_0 \circ a_0)^*\mathcal{C}_{U_0/U} \ar[u]
}
$$
に入る。ここで垂直射は More on Morphisms of Spaces, Lemma
\ref{spaces-more-morphisms-lemma-deform} により同型である。
したがって補題が成り立つ。
\end{proof}

\noindent
群 $\Gamma_0$ を同定しよう。More on Morphisms of Spaces, Remarks
\ref{spaces-more-morphisms-remark-action-by-derivations} および
\ref{spaces-more-morphisms-remark-another-special-case} の議論を図式
$$
\xymatrix{
(U_0 \subset U) \ar@{..>}[rr]_{(e_0, \delta)}
\ar[rd]_{(\text{id}_{U_0}, \text{id}_U)} & &
(R_0 \subset R) \ar[ld]^{(s_0, s)} \\
& (U_0 \subset U)
}
$$
に適用すると、$\delta = \theta \cdot e$ であることが分かる。ここで
$\mathcal{O}_{U_0}$-線形写像
$\theta : e_0^*\Omega_{R_0/U_0} \to \mathcal{C}_{U_0/U}$ は唯一である。したがって、More on Morphisms of Spaces, Lemma
\ref{spaces-more-morphisms-lemma-action-sheaf} を適用することにより全単射
\begin{equation}
\label{spaces-more-groupoids-equation-isomorphism}
\Hom_{\mathcal{O}_{U_0}}(e_0^*\Omega_{R_0/U_0}, \mathcal{C}_{U_0/U})
\longrightarrow
\Gamma_0
\end{equation}
を得る。

\begin{lemma}
\label{spaces-more-groupoids-lemma-composition-is-addition}
全単射 (\ref{spaces-more-groupoids-equation-isomorphism}) は群の同型である。
\end{lemma}

\begin{proof}
$\delta_1, \delta_2 \in \Gamma_0$ が上のように $\theta_1, \theta_2$ に対応し、
合成 $\delta = \delta_1 \circ \delta_2$ で $\Gamma_0$ におけるものが $\theta$ に対応するとする。
$\theta = \theta_1 + \theta_2$ を示さなければならない。More on Morphisms of Spaces, Lemma
\ref{spaces-more-morphisms-lemma-action-by-derivations} を思い出すと、
$\theta_1, \theta_2, \theta$ は導分
$D_1, D_2, D : e_0^{-1}\mathcal{O}_{R_0} \to \mathcal{C}_{U_0/U}$ に対応し、これらは
$D_1 = \theta_1 \circ \text{d}_{R_0/U_0}$ などで与えられる。
$D = D_1 + D_2$ を確かめれば十分である。

\medskip\noindent
等式は茎上で確かめてよい。$\overline{u}$ を $U$ の幾何学的点とし、Section
\ref{spaces-more-groupoids-section-local} で導入した局所環 $A, B, C$ を用いる。
射 $\delta_i$ は環準同型 $\delta_i : B \to A$ に対応する。
$K \subset A$ を平方零イデアルで $A/K = \mathcal{O}_{U_0, \overline{u}}$ となるものとする。
言い換えると、$K$ は $\mathcal{C}_{U_0/U}$ の茎で、$\overline{u}$ におけるものである。
$\delta_i \in \Gamma_0$ であることは、正確に $\delta_i(I) \subset K$ であることを意味する。
導分 $D_i$ は単に写像 $\delta_i - e : B \to A$ である。
$B = s(A) \oplus I$ なので、$D_i$ は $I$ への制限によって決まり、それはちょうど
$\delta_i|_I$ で与えられることが分かる。
% P09-ERR-1349
さらに、$D_i$、したがって $\delta_i$ は $I^2$ を零化する。なぜなら
% P09-ERR-1350
$I = \Ker(I)$ だからである。

\medskip\noindent
証明を終えるため、$\delta$ が合成
$$
B \to C =
(B \otimes_{s, A, t} B)^h_{\mathfrak m_B \otimes B + B \otimes \mathfrak m_B}
\to A
$$
に対応することに注意する。ここで第一の射は $c$ であり、第二の射は規則
$b_1 \otimes b_2 \mapsto \delta_2(t(\delta_1(b_1))) \delta_2(b_2)$
によって決まる。これは (\ref{spaces-more-groupoids-equation-composition}) から従う。
Lemma \ref{spaces-more-groupoids-lemma-first-order-structure-c} により、元 $\zeta$ で $I$ に属するものは
$\zeta \otimes 1 + 1 \otimes \zeta$ と高次の項との和に写ることが分かる。
したがって
$$
D(\zeta) = (\delta_2 \circ t)\left(D_1(\zeta)\right) + D_2(\zeta)
$$
を得る。しかし Lemma \ref{spaces-more-groupoids-lemma-idenity-on-conormal} により、$\delta_2 \circ t$ の作用で
$K = \mathcal{C}_{U_0/U, \overline{u}}$ 上のものは恒等写像である。
これで証明が完了する。
\end{proof}





\section{亜群の性質}
\label{spaces-more-groupoids-section-technical-lemma}

\noindent
この節は More on Groupoids, Section
\ref{more-groupoids-section-technical-lemma} の類似である。
読者には、まずそちらの節を読むことを強く勧める。

\medskip\noindent
次の補題は More on Groupoids, Lemma
\ref{more-groupoids-lemma-property-invariant} の類似である。

\begin{lemma}
\label{spaces-more-groupoids-lemma-property-invariant}
$B \to S$ を Section \ref{spaces-more-groupoids-section-notation} のようなものとする。
$(U, R, s, t, c)$ を $B$ 上の代数空間における亜群とする。次のようにおく。
$\tau \in \{fppf, \linebreak[0] \etale, \linebreak[0]
smooth, \linebreak[0] syntomic\}$。
$\mathcal{P}$ を代数空間の射の性質で、ターゲット上 $\tau$-局所的なものとする
（Descent on Spaces, Definition
\ref{spaces-descent-definition-property-morphisms-local}）。
$\{s : R \to U\}$ および $\{t : R \to U\}$ が $\tau$-位相の被覆であると仮定する。
$W \subset U$ を、$s^{-1}(W) \to W$ が性質 $\mathcal{P}$ をもつような最大の開部分空間とする。
このとき $W$ は $R$-不変である
（Groupoids in Spaces, Definition
\ref{spaces-groupoids-definition-invariant-open}）。
\end{lemma}

\begin{proof}
開部分空間 $W \subset U$ の存在と性質は Descent on Spaces, Lemma
\ref{spaces-descent-lemma-largest-open-of-the-base} に記述されている。
Diagram (\ref{spaces-more-groupoids-equation-diagram}) において $W_1 \subset R$ を、射
$\text{pr}_1 : R \times_{s, U, t} R \to R$ が性質 $\mathcal{P}$ をもつような最大の開部分スキームとする。
前述の Descent on Spaces, Lemma
\ref{spaces-descent-lemma-largest-open-of-the-base} と、$\{s : R \to U\}$ および
$\{t : R \to U\}$ が $\tau$-位相の被覆であるという仮定から、所望の等式
$t^{-1}(W) = W_1 = s^{-1}(W)$ が従う。
\end{proof}

\begin{lemma}
\label{spaces-more-groupoids-lemma-property-G-invariant}
$B \to S$ を Section \ref{spaces-more-groupoids-section-notation} のようなものとする。
$(U, R, s, t, c)$ を $B$ 上の代数空間における亜群とする。
$G \to U$ をその固定化群代数空間とする。次のようにおく。
$\tau \in \{fppf, \linebreak[0] \etale, \linebreak[0]
smooth, \linebreak[0] syntomic\}$。
$\mathcal{P}$ を代数空間の射の性質で、ターゲット上 $\tau$-局所的なものとする。
$\{s : R \to U\}$ および $\{t : R \to U\}$ が $\tau$-位相の被覆であると仮定する。
$W \subset U$ を、$G_W \to W$ が性質 $\mathcal{P}$ をもつような最大の開部分空間とする。
このとき $W$ は $R$-不変である（Groupoids in Spaces, Definition
\ref{spaces-groupoids-definition-invariant-open} を参照せよ）。
\end{lemma}

\begin{proof}
開部分空間 $W \subset U$ の存在と性質は Descent on Spaces, Lemma
\ref{spaces-descent-lemma-largest-open-of-the-base} に記述されている。射
$$
G \times_{U, t} R \longrightarrow R \times_{s, U} G, \quad
(g, r) \longmapsto (r, r^{-1} \circ g \circ r)
$$
は $R$ 上の代数空間の同型である（ここで $\circ$ は亜群における合成を表す）。
したがって $s^{-1}(W) = t^{-1}(W)$ である。これは、前述の Descent on Spaces, Lemma
\ref{spaces-descent-lemma-largest-open-of-the-base} で証明された $W$ の性質による。
\end{proof}




\section{ファイバーの比較}
\label{spaces-more-groupoids-section-fibres}

\noindent
この節は More on Groupoids, Section
\ref{more-groupoids-section-fibres} の類似である。
読者には、まずそちらの節を読むことを強く勧める。

\begin{lemma}
\label{spaces-more-groupoids-lemma-two-fibres}
$B \to S$ を Section \ref{spaces-more-groupoids-section-notation} のようなものとする。
$(U, R, s, t, c)$ を $B$ 上の代数空間における亜群とする。
$K$ を体とし、$r, r' : \Spec(K) \to R$ を射で
$t \circ r = t \circ r' : \Spec(K) \to U$ を満たすものとする。
$u = s \circ r$、$u' = s \circ r'$ とおき、ファイバー積を
$F_u = \Spec(K) \times_{u, U, s} R$ および
$F_{u'} = \Spec(K) \times_{u', U, s} R$ と記す。
このとき $F_u \cong F_{u'}$ であり、これは $K$ 上の代数空間としての同型である。
\end{lemma}

\begin{proof}
Diagram (\ref{spaces-more-groupoids-equation-diagram}) の性質と存在を用いる。
射 $\xi : \Spec(K) \to R \times_{s, U, t} R$ で
$\text{pr}_0 \circ \xi = r$ および $c \circ \xi = r'$ を満たすものが存在する。
$\tilde r = \text{pr}_1 \circ \xi : \Spec(K) \to R$ とおく。
Diagram (\ref{spaces-more-groupoids-equation-diagram}) の下側の二つの正方形を見ると、$F_u$ と $F_{u'}$ はいずれも代数空間
$\Spec(K) \times_{\tilde r, R, \text{pr}_1} (R \times_{s, U, t} R)$
と同一視されることが分かる。
\end{proof}

\noindent
実際、補題の状況では、対の射 $s : (R, r) \to (U, u)$ および
$s : (R, r') \to (U, u')$ は $\tau$-位相において局所的に同型である。
ただし、$\{s: R \to U\}$ は $\tau$-被覆であると仮定する。
必要になれば、ここに精密な主張を挿入する。







\section{亜群の制限}
\label{spaces-more-groupoids-section-restricting-groupoids}

\noindent
この節では、制限によって継承される亜群の性質に関する一連の補題を集める。
これらの補題のほとんどは、制限の定義図式
\begin{equation}
\label{spaces-more-groupoids-equation-restriction}
\vcenter{
\xymatrix{
R' \ar[d] \ar[r] \ar@/_3pc/[dd]_{t'} \ar@/^1pc/[rr]^{s'}&
R \times_{s, U} U' \ar[r] \ar[d] &
U' \ar[d]^g \\
U' \times_{U, t} R \ar[d] \ar[r] &
R \ar[r]^s \ar[d]_t &
U \\
U' \ar[r]^g &
U
}
}
\end{equation}
を考察することで証明できる。Groupoids in Spaces, Lemma
\ref{spaces-groupoids-lemma-restrict-groupoid} を参照せよ。

\begin{lemma}
\label{spaces-more-groupoids-lemma-restrict-preserves-type}
$S$ をスキームとする。$B$ を $S$ 上の代数空間とする。
$(U, R, s, t, c)$ を $B$ 上の代数空間における亜群とする。
$g : U' \to U$ を $B$ 上の代数空間の射とする。
$(U', R', s', t', c')$ を、この $(U, R, s, t, c)$ の $g$ による制限とする。
\begin{enumerate}
\item $s, t$ が局所有限型で $g$ が局所有限型ならば、$s', t'$ は局所有限型である。
\item $s, t$ が局所有限表示で $g$ が局所有限表示ならば、$s', t'$ は局所有限表示である。
\item $s, t$ が平坦で $g$ が平坦ならば、$s', t'$ は平坦である。
% P09-ERR-1351
\item ここにさらに追加する。
\end{enumerate}
\end{lemma}

\begin{proof}
局所有限型であるという性質は合成と任意の基底変換で保たれる。Morphisms of Spaces, Lemmas
\ref{spaces-morphisms-lemma-composition-finite-type} および
\ref{spaces-morphisms-lemma-base-change-finite-type} を参照せよ。
したがって (1) は Diagram (\ref{spaces-more-groupoids-equation-restriction}) から明らかである。
他の場合については Morphisms of Spaces, Lemmas
\ref{spaces-morphisms-lemma-composition-finite-presentation}、
\ref{spaces-morphisms-lemma-base-change-finite-presentation}、
\ref{spaces-morphisms-lemma-composition-flat}、および
\ref{spaces-morphisms-lemma-base-change-flat} を参照せよ。
\end{proof}









\section{体上の群と体上の亜群の性質}
\label{spaces-more-groupoids-section-properties-groupoids-on-fields}

\noindent
読者には、まず亜群スキームに対応する節を見ることを勧める。Groupoids, Section
\ref{groupoids-section-properties-group-schemes-field} および More on Groupoids, Section
\ref{more-groupoids-section-properties-groupoids-on-fields} を参照せよ。

\begin{situation}
\label{spaces-more-groupoids-situation-group-over-field}
ここで $S$ はスキーム、$k$ は $S$ 上の体、$(G, m)$ は $\Spec(k)$ 上の群代数空間である。
\end{situation}

\begin{situation}
\label{spaces-more-groupoids-situation-groupoid-on-field}
ここで $S$ はスキーム、$B$ は代数空間、$(U, R, s, t, c)$ は $B$ 上の代数空間における亜群であり、
$U = \Spec(k)$ と書けるような体 $k$ が存在する。
\end{situation}

\noindent
Situation \ref{spaces-more-groupoids-situation-group-over-field} では、代数空間における亜群
\begin{equation}
\label{spaces-more-groupoids-equation-groupoid-from-group}
(\Spec(k), G, p, p, m)
\end{equation}
が得られることに注意せよ。ここで $p : G \to \Spec(k)$ は $G$ の構造射である。Groupoids in Spaces, Lemma
\ref{spaces-groupoids-lemma-groupoid-from-action} を参照せよ。
これは Situation \ref{spaces-more-groupoids-situation-groupoid-on-field} のような状況である。
この節の残りでは断りなくこれを用いる。

\begin{lemma}
\label{spaces-more-groupoids-lemma-groupoid-on-field-open-multiplication}
Situation \ref{spaces-more-groupoids-situation-groupoid-on-field} において、合成射
$c : R \times_{s, U, t} R \to R$ は平坦かつ普遍的開である。
Situation \ref{spaces-more-groupoids-situation-group-over-field} において、群法則
$m : G \times_k G \to G$ は平坦かつ普遍的開である。
\end{lemma}

\begin{proof}
Diagram (\ref{spaces-more-groupoids-equation-pull}) により、合成は射影写像
$\text{pr}_1 : R \times_{t, U, t} R \to R$ と同型である。
射影は平坦射 $t$ の基底変換なので平坦であり、Morphisms of Spaces, Lemma
\ref{spaces-morphisms-lemma-space-over-field-universally-open} により開である。
第二の主張は第一の主張から直ちに従う。実際、(\ref{spaces-more-groupoids-equation-groupoid-from-group}) では
$m$ が $c$ に一致する。
\end{proof}

\noindent
次の補題は、とりわけ準分離的または局所分離的な代数空間を扱う場合に適用できることに注意せよ
（Decent Spaces, Lemma \ref{decent-spaces-lemma-locally-separated-decent}）。

\begin{lemma}
\label{spaces-more-groupoids-lemma-group-scheme-over-field-separated}
Situation \ref{spaces-more-groupoids-situation-groupoid-on-field} において $R$ が適正な空間であると仮定する。
このとき $R$ は分離的代数空間である。
Situation \ref{spaces-more-groupoids-situation-group-over-field} において $G$ が適正な代数空間であると仮定する。
% P09-ERR-1352
このとき $G$ は分離的代数空間である。
\end{lemma}

\begin{proof}
まず第二の主張を証明する。Groupoids in Spaces, Lemma
\ref{spaces-groupoids-lemma-group-scheme-separated} により、
% P09-ERR-1353
$e : S \to G$ が閉埋め込みであることを示さなければならない。
これは Decent Spaces, Lemma
\ref{decent-spaces-lemma-finite-residue-field-extension-finite} から従う。

\medskip\noindent
次に第一の主張を証明する。このため $B$ を $S$ に置き換えてよい。
% P09-ERR-1354
上の段落により固定化群スキーム $G \to U$ は分離的である。
Groupoids in Spaces, Lemma \ref{spaces-groupoids-lemma-diagonal} により、射
$j = (t, s) : R \to U \times_S U$ は分離的である。
$U$ は体のスペクトルなので、スキーム $U \times_S U$ はアフィンである
（Schemes, Section \ref{schemes-section-fibre-products} におけるファイバー積の構成による）。
したがって $R$ は分離的である。Morphisms of Spaces, Lemma
\ref{spaces-morphisms-lemma-separated-over-separated} を参照せよ。
\end{proof}

\begin{lemma}
\label{spaces-more-groupoids-lemma-restrict-groupoid-on-field}
% P09-ERR-1355
Situation \ref{spaces-more-groupoids-situation-groupoid-on-field} において考える。
$k'/k$ を体の拡大、$U' = \Spec(k')$ とし、$(U', R', s', t', c')$ を
$(U, R, s, t, c)$ の $U' \to U$ による制限とする。定義図式
$$
\xymatrix{
R' \ar[d] \ar[r] \ar@/_3pc/[dd]_{t'} \ar@/^1pc/[rr]^{s'} \ar@{..>}[rd] &
R \times_{s, U} U' \ar[r] \ar[d] &
U' \ar[d] \\
U' \times_{U, t} R \ar[d] \ar[r] &
R \ar[r]^s \ar[d]_t &
U \\
U' \ar[r] &
U
}
$$
におけるすべての射は全射、平坦、かつ普遍的開である。
点線の射 $R' \to R$ は、さらにアフィンである。
\end{lemma}

\begin{proof}
射 $U' \to U$ は $\Spec(k') \to \Spec(k)$ に等しく、したがってアフィン、全射、かつ平坦である。
射 $s, t : R \to U$ および射 $U' \to U$ は Morphisms, Lemma
% P09-ERR-1356
\ref{morphisms-lemma-scheme-over-field-universally-open} により普遍的開である。
$R$ は空でなく $U$ は体のスペクトルなので、射 $s, t : R \to U$ は全射かつ平坦である。
このとき、Morphisms of Spaces, Lemmas
\ref{spaces-morphisms-lemma-base-change-surjective}、
\ref{spaces-morphisms-lemma-composition-surjective}、
\ref{spaces-morphisms-lemma-composition-open}、
\ref{spaces-morphisms-lemma-base-change-affine}、
\ref{spaces-morphisms-lemma-composition-affine}、
\ref{spaces-morphisms-lemma-base-change-flat}、および
\ref{spaces-morphisms-lemma-composition-flat} を用いて結論を得る。
\end{proof}

\begin{lemma}
\label{spaces-more-groupoids-lemma-groupoid-on-field-explain-points}
Situation \ref{spaces-more-groupoids-situation-groupoid-on-field} において考える。
任意の点 $r \in |R|$ に対して、次のものが存在する。
\begin{enumerate}
\item 体拡大 $k'/k$。ここで $k'$ は代数閉体である。
\item 点 $r' : \Spec(k') \to R'$。ここで
$(U', R', s', t', c')$ は $(U, R, s, t, c)$ の
$\Spec(k') \to \Spec(k)$ による制限である。
\end{enumerate}
これらは次を満たす。
\begin{enumerate}
\item 点 $r'$ の像は $r$ である（射 $R' \to R$ のもとで）。
\item 写像
$s' \circ r', t' \circ r' : \Spec(k') \to \Spec(k')$
は自己同型である。
\end{enumerate}
\end{lemma}

\begin{proof}
点 $r$ を射 $r : \Spec(K) \to R$ によって表す。ここで $K$ はある体である。
補題を証明するには、代数閉体 $k'$ と可換図式
$$
\xymatrix{
k' & k' \ar[l]^1 & \\
k' \ar[u]^\tau & K \ar[lu]^\sigma & k \ar[l]^-s \ar[lu]_i \\
& k \ar[lu]^i \ar[u]_t
}
$$
を見いだせばよい。ここで $s, t : k \to K$ は、
$s \circ r$ および $t \circ r$ から得られる体準同型である。
More on Groupoids, Lemma
\ref{more-groupoids-lemma-groupoid-on-field-explain-points}
の証明には、このような図式の構成法が示されている。
\end{proof}

\begin{lemma}
\label{spaces-more-groupoids-lemma-groupoid-on-field-move-point}
Situation \ref{spaces-more-groupoids-situation-groupoid-on-field} において考える。
$r : \Spec(k) \to R$ が、$s \circ r, t \circ r$ が
$\Spec(k)$ の自己同型となるような射ならば、写像
$$
R \longrightarrow R, \quad
x \longmapsto c(r, x)
$$
は自己同型 $R \to R$ であり、$e$ を $r$ に写す。
\end{lemma}

\begin{proof}
証明は More on Groupoids, Lemma
\ref{more-groupoids-lemma-groupoid-on-field-move-point}
の証明と同一である。
\end{proof}

\begin{lemma}
% P09-ERR-1357
\label{spaces-more-groupoids-lemma-groupoid-on-field-geometrically-irreducible}
Situation \ref{spaces-more-groupoids-situation-groupoid-on-field} において、代数空間 $R$ は幾何学的単枝である。
Situation \ref{spaces-more-groupoids-situation-group-over-field} において、代数空間 $G$ は幾何学的単枝である。
\end{lemma}

\begin{proof}
$r \in |R|$ とする。$R$ が $r$ において幾何学的単枝であることを示さなければならない。
Lemma \ref{spaces-more-groupoids-lemma-restrict-groupoid-on-field} と
Descent on Spaces, Lemma \ref{spaces-descent-lemma-descend-unibranch}
を組み合わせると、$k$ が代数閉で、$r$ が
射 $r : \Spec(k) \to R$ から来て、$s \circ r$ と $t \circ r$ が
$\Spec(k)$ の自己同型となる場合を証明すれば十分である。
Lemma \ref{spaces-more-groupoids-lemma-groupoid-on-field-move-point} により、
$r = e$ が $R$ の単位元であり、$k$ が代数閉である場合に帰着する。

\medskip\noindent
$r = e$ かつ $k$ が代数閉であると仮定する。
$A = \mathcal{O}_{R, e}$ を $R$ の $e$ におけるエタール局所環とし、
$C = \mathcal{O}_{R \times_{s, U, t} R, (e, e)}$ を
$R \times_{s, U, t} R$ の $(e, e)$ におけるエタール局所環とする。
More on Algebra, Lemma
\ref{more-algebra-lemma-minimal-primes-tensor-strictly-henselian}
により、$\mathfrak q$ が $C$ の極小素イデアルであることと、
$1$ 対 $1$ に対応して $\mathfrak p, \mathfrak p' \subset A$ が
極小素イデアルの対であることとは同値である。
一方、合成則は平坦な環準同型
$$
\xymatrix{
A \ar[r]_{c^\sharp} & C & \mathfrak q \\
& A \otimes_{s^\sharp, k, t^\sharp} A \ar[u] &
\mathfrak p \otimes A + A \otimes \mathfrak p' \ar@{|}[u]
}
$$
を誘導する。$(c^\sharp)^{-1}(\mathfrak q)$ は
$\mathfrak p$ と $\mathfrak p'$ の双方を含むことに注意せよ。これは図式
$$
\xymatrix{
A \ar[r]_{c^\sharp} & C \\
A \otimes_{s^\sharp, k} k \ar[u] &
A \otimes_{s^\sharp, k, t^\sharp} A \ar[l]_{1 \otimes e^\sharp} \ar[u]
}
\quad\quad
\xymatrix{
A \ar[r]_{c^\sharp} & C \\
k \otimes_{k, t^\sharp} A \ar[u] &
A \otimes_{s^\sharp, k, t^\sharp} A \ar[l]_{e^\sharp \otimes 1} \ar[u]
}
$$
が (\ref{spaces-more-groupoids-equation-diagram}) により可換だからである。
$c^\sharp$ は平坦であるから（$c$ は
Lemma \ref{spaces-more-groupoids-lemma-groupoid-on-field-open-multiplication} により平坦射である）、
$(c^\sharp)^{-1}(\mathfrak q)$ は $A$ の極小素イデアルである。
したがって
$\mathfrak p = (c^\sharp)^{-1}(\mathfrak q) = \mathfrak p'$ である。
\end{proof}

\noindent
次の補題では、Properties of Spaces, Section
\ref{spaces-properties-section-dimension} で定義された代数空間の（点における）次元を用いる。
また、Properties of Spaces, Section
\ref{spaces-properties-section-dimension-local-ring} で定義された局所環の次元と、
Morphisms of Spaces, Section \ref{spaces-morphisms-section-relative-dimension}
にある点の超越次数も用いる。

\begin{lemma}
\label{spaces-more-groupoids-lemma-groupoid-on-field-locally-finite-type-dimension}
Situation \ref{spaces-more-groupoids-situation-groupoid-on-field} において、$s, t$ が局所有限型であると仮定する。
すべての $r \in |R|$ に対して、次が成り立つ。
\begin{enumerate}
\item $\dim(R) = \dim_r(R)$、
\item $r$ の $\Spec(k)$ 上での $s$ を介した超越次数は、
$r$ の $\Spec(k)$ 上での $t$ を介した超越次数に等しく、
\item (2) で述べた超越次数が $0$ ならば、
$\dim(R) = \dim(\mathcal{O}_{R, \overline{r}})$ である。
\end{enumerate}
\end{lemma}

\begin{proof}
$r \in |R|$ とする。$\text{trdeg}(r/_{\!\! s}k)$ は、
$r$ の $\Spec(k)$ 上での $s$ を介した超越次数を表すものとする。
エタール射 $\varphi : V \to R$ を選ぶ。ここで $V$ はスキームで、
$v \in V$ は $r$ に写る。補題の前で述べた定義を用いると、
$$
\dim_r(R) = \dim_v(V) =
\dim(\mathcal{O}_{V, v}) + \text{trdeg}_{s(k)}(\kappa(v)) =
\dim(\mathcal{O}_{R, \overline{r}}) + \text{trdeg}(r/_{\!\! s}k)
$$
となり、$t$ についても同様である（第二の等号は Morphisms, Lemma
\ref{morphisms-lemma-dimension-fibre-at-a-point} による）。したがって
$\text{trdeg}(r/_{\!\! s}k) = \text{trdeg}(r/_{\!\! t}k)$
であり、(2) が成り立つ。

\medskip\noindent
$k'/k$ を体拡大とする。$R'$ を $R$ の $\Spec(k')$ への制限とする
（Lemma \ref{spaces-more-groupoids-lemma-restrict-groupoid-on-field} を参照）、
体の射による二回の基底変換によって $R$ から得られることに注意せよ。
したがって Morphisms of Spaces, Lemma
\ref{spaces-morphisms-lemma-dimension-fibre-after-base-change}
により、この操作では点における $R$ の次元は変わらない。
ゆえに (1) を証明するには、Lemma
\ref{spaces-more-groupoids-lemma-groupoid-on-field-explain-points} により、$r$ が射
$r : \Spec(k) \to R$ で表され、$s \circ r$ と $t \circ r$ の双方が
$\Spec(k)$ の自己同型であると仮定してよい。
この場合、自己同型 $R \to R$ が存在し、$r$ を $e$ に写す
（Lemma \ref{spaces-more-groupoids-lemma-groupoid-on-field-move-point}）。
したがって $\dim_r(R) = \dim_e(R)$ は任意の $r$ に対して成り立つ。
定義により、これは $\dim_r(R) = \dim(R)$ を意味する。

\medskip\noindent
(3) は、上の議論で得た結果から形式的に従う。
\end{proof}

\begin{lemma}
\label{spaces-more-groupoids-lemma-group-over-field-locally-finite-type-dimension}
Situation \ref{spaces-more-groupoids-situation-group-over-field} において、
% P09-ERR-1358
$G$ が局所有限型であると仮定する。
すべての $g \in |G|$ に対して、次が成り立つ。
\begin{enumerate}
\item $\dim(G) = \dim_g(G)$、
\item $g$ の $k$ 上の超越次数が $0$ ならば、
$\dim(G) = \dim(\mathcal{O}_{G, \overline{g}})$ である。
\end{enumerate}
\end{lemma}

\begin{proof}
Lemma \ref{spaces-more-groupoids-lemma-groupoid-on-field-locally-finite-type-dimension}
を (\ref{spaces-more-groupoids-equation-groupoid-from-group}) に適用すれば直ちに従う。
\end{proof}

\begin{lemma}
\label{spaces-more-groupoids-lemma-groupoid-on-field-dimension-equal-stabilizer}
Situation \ref{spaces-more-groupoids-situation-groupoid-on-field} において、$s, t$ が局所有限型であると仮定する。
$G = \Spec(k)\allowbreak
\times_{\Delta, \Spec(k) \times_B \Spec(k), t \times s} R$
を固定化群代数空間とする。このとき $\dim(R) = \dim(G)$ である。
\end{lemma}

\begin{proof}
$G$ と $R$ は等次元なので（Lemmas
\ref{spaces-more-groupoids-lemma-groupoid-on-field-locally-finite-type-dimension} および
\ref{spaces-more-groupoids-lemma-group-over-field-locally-finite-type-dimension} を参照）、
$\dim_e(R) = \dim_e(G)$ を証明すれば十分である。
$V$ をアフィンスキーム、$v \in V$ とし、
% P09-ERR-1359
$\varphi : V \to R$ を $\varphi(v) = e$ を満たすスキームのエタール射とする。
$V$ は Noether スキームであることに注意せよ。実際、$s \circ \varphi$ は
局所有限型射の合成として局所有限型であり、かつ $V$ は準コンパクトである
（Morphisms of Spaces, Lemmas
\ref{spaces-morphisms-lemma-composition-finite-type}、
\ref{spaces-morphisms-lemma-etale-locally-finite-presentation}、
\ref{spaces-morphisms-lemma-finite-presentation-finite-type}、および
Morphisms, Lemma \ref{morphisms-lemma-finite-type-noetherian} を用いる）。
したがって $V$ は局所連結である（Properties, Lemma
\ref{properties-lemma-Noetherian-topology} および
Topology, Lemma \ref{topology-lemma-locally-Noetherian-locally-connected} を参照）。
よって $V$ を $v$ を含む連結成分で置き換えてよい
（それは $V$ の開かつ閉な部分スキームなので、依然としてアフィンである）。
$T = V_{red}$ を $V$ の被約化とする。二つの射
$a, b : T \to \Spec(k)$ を考える。これらは
$a = s \circ \varphi|_T$ および $b = t \circ \varphi|_T$ で与えられる。
$a, b$ は同じ体準同型 $k \to \kappa(v)$ を誘導することに注意せよ。
これは $\varphi(v) = e$ だからである。
$k_a \subset \Gamma(T, \mathcal{O}_T)$ を
$a^\sharp(k) \subset \Gamma(T, \mathcal{O}_T)$ の整閉包とする。同様に、
$k_b \subset \Gamma(T, \mathcal{O}_T)$ を
$b^\sharp(k) \subset \Gamma(T, \mathcal{O}_T)$ の整閉包とする。
Varieties, Proposition \ref{varieties-proposition-unique-base-field}
により $k_a = k_b$ である。したがって、次の可換図式を得る。
$$
\xymatrix{
k \ar[rd]^a \ar[rrrd] \\
& k_a = k_b \ar[r] & \Gamma(T, \mathcal{O}_T) \ar[r] & \kappa(v) \\
k \ar[ru]_b \ar[rrru]
}
$$
上で論じたように、長い矢印は等しい。
$k_a = k_b \to \kappa(v)$ は単射なので、二つの射 $a$ と $b$ は一致する。
したがって $T \to R$ は $G$ を経由する。
ゆえに $R_{red} = G_{red}$ は $e$ のある開近傍で成り立ち、これは確かに
$\dim_e(R) = \dim_e(G)$ を含意する。
\end{proof}

\section{体上の群代数空間}
\label{spaces-more-groupoids-section-group-algebraic-spaces-over-fields}

\noindent
体上には非分離的な群代数空間が存在する。実際、標数零の体上の
$\mathbf{G}_a/\mathbf{Z}$ がその例である。Examples, Section
\ref{examples-section-non-separated-group-space} を参照せよ。
任意の体上の群スキームは分離的であるから
（Lemma \ref{spaces-more-groupoids-lemma-group-scheme-over-field-separated}）、
体上の非分離的な群代数空間はいずれも表現可能ではない。
一方、体上の群代数空間は、適正でさえあれば分離的である。
Lemma \ref{spaces-more-groupoids-lemma-group-scheme-over-field-separated} を参照せよ。
この節では、体上の分離的な群代数空間が表現可能、すなわちスキームであることを示す。

\begin{lemma}
\label{spaces-more-groupoids-lemma-group-space-scheme-over-kbar}
$k$ を代数閉包 $\overline{k}$ をもつ体とする。
$G$ を $k$ 上の分離的な群代数空間とする\footnote{$G$ が適正であると仮定すれば十分である。
たとえば、Lemma \ref{spaces-more-groupoids-lemma-group-scheme-over-field-separated} により、
局所分離的または準分離的であればよい。}。
このとき $G_{\overline{k}}$ はスキームである。
\end{lemma}

\begin{proof}
Spaces over Fields, Lemma
\ref{spaces-over-fields-lemma-when-scheme-after-base-change}
により、$G_K$ がスキームとなるような体拡大 $K/k$ が存在することを示せば十分である。
$G_K' \subset G_K$ を、Properties of Spaces, Lemma
\ref{spaces-properties-lemma-subscheme} における $G_K$ のスキーム的軌跡と記す。
Properties of Spaces, Proposition
\ref{spaces-properties-proposition-locally-quasi-separated-open-dense-scheme}
により $G_K' \subset G_K$ は稠密開であり、特に空でない。
スキーム $U$ と全射エタール射 $U \to G$ を選ぶ。
Varieties, Lemma \ref{varieties-lemma-make-Jacobson} により、
$K$ が十分大きな超越次数をもつ代数閉体ならば、$U_K$ は Jacobson スキームであり、
$U_K$ のすべての閉点は $K$-有理的である。
したがって $G_K'$ は $K$-有理点をもつ。また、任意の $K$-有理点を
$G_K$ から取るとき、それが $G_K'$ に属することを示せば十分である。
$g \in G_K(K)$ を $K$-有理点、$g' \in G_K'(K)$ をスキーム的軌跡内の
$K$-有理点とする。このとき、$g$ は $G_K'$ の、自己同型
$$
G_K \longrightarrow G_K,\quad
h \longmapsto g(g')^{-1}h
$$
という $G_K$ の自己同型による像に属する。代数空間としての $G_K$ の自己同型は
$G_K'$ を保つので、所望のとおり $g \in G_K'$ である。
\end{proof}

\begin{lemma}
\label{spaces-more-groupoids-lemma-group-space-scheme-locally-finite-type-over-k}
$k$ を体とし、$G$ を $k$ 上の群代数空間とする。
$G$ が $k$ 上分離的かつ局所有限型ならば、$G$ はスキームである。
\end{lemma}

\begin{proof}
これは Lemma \ref{spaces-more-groupoids-lemma-group-space-scheme-over-kbar}、
Groupoids, Lemma \ref{groupoids-lemma-points-in-affine}、および
Spaces over Fields, Lemma
\ref{spaces-over-fields-lemma-scheme-over-algebraic-closure-enough-affines}
から従う。
\end{proof}

\begin{proposition}
\label{spaces-more-groupoids-proposition-group-space-scheme-over-field}
$k$ を体とし、$G$ を $k$ 上の群代数空間とする。
$G$ が分離的ならば、$G$ はスキームである。
\end{proposition}

\begin{proof}
% P09-ERR-1360
この補題は Lemma \ref{spaces-more-groupoids-lemma-group-space-scheme-locally-finite-type-over-k}
（これは実際に関心をもつすべての場合を扱う）を一般化する。
証明は Spaces over Fields, Lemma
\ref{spaces-over-fields-lemma-scheme-over-algebraic-closure-enough-affines}
の証明と非常によく似ている。この証明は Lemma
\ref{spaces-more-groupoids-lemma-group-space-scheme-locally-finite-type-over-k}
の証明でも用いた。まずそちらの証明を読むことを勧める。

\medskip\noindent
Lemma \ref{spaces-more-groupoids-lemma-group-space-scheme-over-kbar} により、基底変換
$G_{\overline{k}}$ はスキームである。
$K/k$ を超越次数が非常に大きい純超越拡大とする。
Spaces over Fields, Lemma
\ref{spaces-over-fields-lemma-scheme-after-purely-transcendental-base-change}
により、$G_K$ がスキームであることを示せば十分である。
$K^{perf}$ を $K$ の完全閉包とする。Spaces over Fields, Lemma
\ref{spaces-over-fields-lemma-scheme-after-purely-inseparable-base-change}
により、$G_{K^{perf}}$ がスキームであることを示せば十分である。
% P09-ERR-1361
包含 $K \subset K^{perf} \subset \overline{K}$ をとり、これを $K$ の代数閉包とする。
埋め込み $\overline{k} \to \overline{K}$ を $k$ 上で選べるので、
$G_{\overline{K}}$ はスキーム $G_{\overline{k}}$ の
$\overline{k} \to \overline{K}$ による基底変換である。
Varieties, Lemma \ref{varieties-lemma-make-Jacobson} により、
$G_{\overline{K}}$ は、すべての閉点の剰余体が $\overline{K}$ である Jacobson スキームである。

\medskip\noindent
$G_{\overline{K}} \to G_{K^{perf}}$ は全射なので、
$g \in |G_{K^{perf}}|$ をその像とする。これが $G_{\overline{K}}$ の任意の閉点の像であるとき、
% P09-ERR-1362
$G_K$ のスキーム的軌跡に属することを示せば十分である。
特に、$g$ を射 $g : \Spec(L) \to G_{K^{perf}}$ で表せる。
ここで $L/K^{perf}$ は分離代数拡大である
（たとえば $L = \overline{K}$ ととれる）。したがってスキーム
\begin{align*}
T & = \Spec(L) \times_{G_{K^{perf}}} G_{\overline{K}} \\
& =
\Spec(L) \times_{\Spec(K^{perf})} \Spec(\overline{K}) \\
& =
\Spec(L \otimes_{K^{perf}} \overline{K})
\end{align*}
は、$\overline{K}$ のコピーの有限積である代数のフィルター付き余極限となる
$\overline{K}$-代数のスペクトルである。
したがって Groupoids, Lemma \ref{groupoids-lemma-compact-set-in-affine}
により、アフィン開 $W \subset G_{\overline{K}}$ を見いだせる。これは
$g_{\overline{K}} : T \to G_{\overline{K}}$ の像を含む。

\medskip\noindent
準コンパクト開 $V \subset G_{K^{perf}}$ を選び、これが $W$ の像を含むようにする。
Spaces over Fields, Lemma
\ref{spaces-over-fields-lemma-when-scheme-after-base-change}
により、$V_{K'}$ がスキームとなるような有限拡大 $K'/K^{perf}$ が存在する。
$K'$ を拡大すれば、アフィン開 $U' \subset V_{K'} \subset G_{K'}$ が存在し、
その $\overline{K}$ への基底変換が $W$ を復元すると仮定してよい
（$V_{\overline{K}}$ はスキーム $V_{K''}$ の極限であり、ここで
$K' \subset K'' \subset \overline{K}$ は有限であることと、Limits, Lemmas
\ref{limits-lemma-descend-opens} および
\ref{limits-lemma-limit-affine} を用いる）。
% P09-ERR-1363
$K'/K^{perf}$ は Galois 拡大であると仮定してよい
（正規閉包をとり、Fields, Lemma \ref{fields-lemma-normal-closure} と
$K^{perf}$ が完全であることを用いる）。$H = \text{Gal}(K'/K^{perf})$ とおく。
構成により、$H$-不変な閉部分スキーム
$\Spec(L) \times_{G_{K^{perf}}} G_{K'}$ は $U'$ に含まれる。
Spaces over Fields, Lemmas
\ref{spaces-over-fields-lemma-base-change-by-Galois} および
\ref{spaces-over-fields-lemma-when-quotient-scheme-at-point}
により結論を得る。
\end{proof}





\section{群上の有理曲線の非存在}
\label{spaces-more-groupoids-section-no-rational-curves}

\noindent
この節では、$\mathbf{P}^1$ から体上局所有限型の群代数空間への
非定数射が存在しないことを証明する。

\begin{lemma}
\label{spaces-more-groupoids-lemma-factor-through-over-open}
$S$ をスキーム、$B$ を $S$ 上の代数空間とする。
$f : X \to Y$ および $g : X \to Z$ を $B$ 上の代数空間の射とする。
次を仮定する。
\begin{enumerate}
\item $Y \to B$ は分離的である。
\item $g$ は全射、平坦、かつ局所有限表示である。
\item スキーム論的に稠密な開 $V \subset Z$ が存在し、
$f|_{g^{-1}(V)} : g^{-1}(V) \to Y$ は $V$ を経由する。
\end{enumerate}
このとき $f$ は $g$ を経由する。
\end{lemma}

\begin{proof}
$R = X \times_Z X$ とおく。(2) により層として $Z = X/R$ である。
また (2) により、$V$ の $R$ における逆像は $R$ でスキーム論的に稠密である
（Morphisms of Spaces, Lemma
\ref{spaces-morphisms-lemma-flat-morphism-scheme-theoretically-dense-open}）。
% P09-ERR-1364
このとき二つの合成 $R \to X \to Y$ は Morphisms of Spaces, Lemma
\ref{spaces-morphisms-lemma-equality-of-morphisms} により等しい。
補題が従う。
\end{proof}

\begin{lemma}
\label{spaces-more-groupoids-lemma-quotient-power-P1}
\begin{slogan}
体上の射影直線の空でない積から、体上の分離的有限型代数空間への射は、
射影直線の積への射影の後に有限射として因子分解する。
\end{slogan}
$k$ を体とする。$n \geq 1$ とし、$(\mathbf{P}^1_k)^n$ を
$n$ 重の $\Spec(k)$ 上の自己積とする。
$f : (\mathbf{P}^1_k)^n \to Z$ を $k$ 上の代数空間の射とする。
$Z$ が $k$ 上分離的有限型ならば、$f$ は次のように因子分解する。
$$
(\mathbf{P}^1_k)^n \xrightarrow{projection}
(\mathbf{P}^1_k)^m \xrightarrow{finite} Z.
$$
\end{lemma}

\begin{proof}
$k$ は代数閉であると仮定してよい（詳細は省略する）。これは有理点を用いて議論するためだけであり、
望むならこの仮定を回避できる。証明中の積は $k$ 上でとる。
$(\mathbf{P}^1_k)^n$ の自己同型群代数空間は
$G = (\text{GL}_{2, k})^n$ を含む。
$C \subset (\mathbf{P}^1_k)^n$ を一点に写される閉部分多様体
（特に $k$ 上既約）とすると、More on Morphisms of Spaces, Lemma
\ref{spaces-more-morphisms-lemma-flat-proper-family-cannot-collapse-fibre}
を射
$$
G \times C \to G \times Z,\quad (g, c) \mapsto (g, f(g \cdot c))
$$
に $G$ 上で適用できる。したがって $g(C)$ は一点に写される。
ここで $g \in G(k)$ は Zariski 開 $U \subset G$ に属する。
$x = (x_1, \ldots, x_n)$、$y = (y_1, \ldots, y_n)$ を
$k$-値点とし、これらは $(\mathbf{P}^1_k)^n$ の点である。
$I \subset \{1, \ldots, n\}$ を添字 $i$ の集合とし、
その各添字は $x_i = y_i$ を満たすものとする。このとき
$$
\{g(x) \mid g(y) = y,\ g \in U(k)\}
$$
は射影 $\pi_I : (\mathbf{P}^1_k)^n \to \prod_{i \in I} \mathbf{P}^1_k$
のファイバーで Zariski 稠密である（演習）。
したがって、相異なる $x, y \in C(k)$ に対し、$f$ は
$\pi_I$ のファイバー全体を一点に写す。このファイバーは $x, y$ を含む。
さらに、$U(k)$-軌道による $C$ の像は $\pi_I$ のファイバーの Zariski 開集合と交わる。
Lemma \ref{spaces-more-groupoids-lemma-factor-through-over-open} により、射 $f$ は $\pi_I$ を経由する。
この過程を有限回繰り返すと、$f$ のすべてのファイバーが $k$-点上で有限である段階に達する。
この場合 $f$ は有限である。これは More on Morphisms of Spaces, Lemma
\ref{spaces-more-morphisms-lemma-proper-finite-fibre-finite-in-neighbourhood}
と、$k$-点が $Z$ で稠密であること
（Spaces over Fields, Lemma
\ref{spaces-over-fields-lemma-smooth-separable-closed-points-dense}）による。
\end{proof}

\begin{lemma}
\label{spaces-more-groupoids-lemma-no-nonconstant-morphism-from-P1-to-group}
\begin{slogan}
群上には完備有理曲線は存在しない。
\end{slogan}
$k$ を体とし、$G$ を $k$ 上局所有限型の分離的群代数空間とする。
非定数射 $f : \mathbf{P}^1_k \to G$ は $\Spec(k)$ 上には存在しない。
\end{lemma}

\begin{proof}
$f$ が非定数であると仮定する。射
% P09-ERR-1365
$$
\mathbf{P}^1_k \times_{\Spec(k)} \ldots \times_{\Spec(k)} \mathbf{P}^1_k
\longrightarrow G,
\quad (t_1, \ldots, t_n) \longmapsto f(g_1) \ldots f(g_n)
$$
を考える。右辺では群の積を用いている。
Lemma \ref{spaces-more-groupoids-lemma-quotient-power-P1} と $f$ が非定数であるという仮定により、
この射はその像上有限である。したがって $\dim(G) \geq n$ は
すべての $n$ に対して成り立つが、これは Lemma
\ref{spaces-more-groupoids-lemma-group-over-field-locally-finite-type-dimension} と、
$G$ が $k$ 上局所有限型であることに反する。
\end{proof}

\section{射の有限部分}
\label{spaces-more-groupoids-section-finite}

\noindent
$S$ をスキームとする。
$f : X \to Y$ を $S$ 上の代数空間の射とする。
代数空間またはスキーム $T$ をとり、これが $S$ 上にあるとする。
次を満たす対 $(a, Z)$ を考える。
\begin{equation}
\label{spaces-more-groupoids-equation-finite-conditions}
\begin{matrix}
a : T \to Y\text{ is a morphism over }S, \\
Z \subset T \times_Y X\text{ is an open subspace} \\
\text{such that }\text{pr}_0|_Z : Z \to T\text{ is finite.}
\end{matrix}
\end{equation}
$h : T' \to T$ を $S$ 上の代数空間の射とし、$(a, Z)$ を
(\ref{spaces-more-groupoids-equation-finite-conditions}) にある $T$ 上の対とする。
$a' = a \circ h$ および
$Z' = (h \times \text{id}_X)^{-1}(Z) = T' \times_T Z$ とおく。
このとき $(a', Z')$ は (\ref{spaces-more-groupoids-equation-finite-conditions}) にある $T'$ 上の対である。
これは有限射が基底変換で保たれることから従う。Morphisms of Spaces, Lemma
\ref{spaces-morphisms-lemma-base-change-integral} を参照せよ。
したがって関手
\begin{equation}
\label{spaces-more-groupoids-equation-finite}
\begin{matrix}
(X/Y)_{fin} : &
(\Sch/S)^{opp} &
\longrightarrow &
\textit{Sets} \\
& T & \longmapsto &
\{(a, Z)\text{ as above}\}
\end{matrix}
\end{equation}
を得る。応用上は主として、この関手 $(X/Y)_{fin}$ において
$f$ が分離的かつ局所有限型である場合に関心がある。その内容を概観するには、
Remark \ref{spaces-more-groupoids-remark-finite-quasi-finite-separated-morphism-schemes} を参照せよ。

\begin{lemma}
\label{spaces-more-groupoids-lemma-finite-sheaf}
$S$ をスキームとする。
$f : X \to Y$ を $S$ 上の代数空間の射とする。このとき次が成り立つ。
\begin{enumerate}
\item 前層 $(X/Y)_{fin}$ は fppf 位相に関する層条件を満たす。
\item $T$ が $S$ 上の代数空間ならば、標準的な全単射
$$
\Mor_{\Sh((\Sch/S)_{fppf})}(T, (X/Y)_{fin})
=
\{(a, Z)\text{ satisfying \ref{spaces-more-groupoids-equation-finite-conditions}}\}
$$
が存在する。
\end{enumerate}
\end{lemma}

\begin{proof}
$T$ を $S$ 上の代数空間とする。
$\{T_i \to T\}$ を（代数空間による）fppf 被覆とする。
$s_i = (a_i, Z_i)$ を \ref{spaces-more-groupoids-equation-finite-conditions} を満たす
$T_i$ 上の対とし、
$s_i|_{T_i \times_T T_j} = s_j|_{T_i \times_T T_j}$ が成り立つと仮定する。
第一に、これは特に $a_i$ と $a_j$ が同じ射
$T_i \times_T T_j \to Y$ を定めることを含意する。
Descent on Spaces, Lemma
\ref{spaces-descent-lemma-fpqc-universal-effective-epimorphisms}
により、一意な射 $a : T \to Y$ が存在し、$a_i$ は合成
$T_i \to T \to Y$ に等しい。
第二に、$Z_i \subset T_i \times_Y X$ は、その
$(T_i \times_T T_j) \times_Y X$ における逆像が等しい開部分空間である。
$\{T_i \times_Y X \to T \times_Y X\}$ は fppf 被覆なので、
一意な開部分空間 $Z \subset T \times_Y X$ が存在し、その制限は
$Z_i$ である（$T_i$ 上で）。Descent on Spaces, Lemma
\ref{spaces-descent-lemma-open-fpqc-covering} を参照せよ。
射影 $Z \to T$ は有限であると主張する。
これは有限性が fpqc 位相について局所的であることから従う。Descent on Spaces, Lemma
\ref{spaces-descent-lemma-descending-property-finite} を参照せよ。

\medskip\noindent
前段落の結果は、特に (1) を含意することに注意せよ。

\medskip\noindent
$T$ を $S$ 上の代数空間とする。(2) を証明するため、表示された集合の間に
互いに逆な写像を構成する。以下「対」と言うときは、条件
\ref{spaces-more-groupoids-equation-finite-conditions} を満たす対を意味する。

\medskip\noindent
$v : T \to (X/Y)_{fin}$ を自然変換とする。
スキーム $U$ と全射エタール射 $p : U \to T$ を選ぶ。
このとき $v(p) \in (X/Y)_{fin}(U)$ は対 $(a_U, Z_U)$ に対応し、この対は $U$ 上にある。
$R = U \times_T U$ とし、$t, s : R \to U$ を射影とする。
$v$ は関手の変換なので、$(a_U, Z_U)$ の $s$ と $t$ による引き戻しは一致する。
したがって $\{U \to T\}$ は fppf 被覆なので、第一段落の結果を適用し、
% P09-ERR-1366
一意な対 $(a, Z)$ が存在し、これは $T$ 上にあると結論できる。

\medskip\noindent
逆に、$(a, Z)$ を $T$ 上の対とする。
$U \to T$、$R = U \times_T U$、および $t, s : R \to U$ を上のとおりとする。
このとき制限 $(a, Z)|_U$ は、Yoneda の補題
（Categories, Lemma \ref{categories-lemma-yoneda}）により関手の変換
$v : h_U \to (X/Y)_{fin}$ を与える。
二つの引き戻し $s^*(a, Z)|_U$ と $t^*(a, Z)|_U$ は等しいので、
$v$ は二つの写像 $h_t, h_s : h_R \to h_U$ を余等化する。
Spaces, Lemma \ref{spaces-lemma-space-presentation} により
$T = U/R$ は fppf 商層であり、(1) により $(X/Y)_{fin}$ は fppf 層なので、
$v$ は写像 $T \to (X/Y)_{fin}$ を経由する。

\medskip\noindent
上の二つの構成が互いに逆であることの確認は省略する。
\end{proof}

\begin{lemma}
\label{spaces-more-groupoids-lemma-finite-open}
$S$ をスキームとする。可換図式
$$
\xymatrix{
X' \ar[rr]_j \ar[rd] & & X \ar[ld] \\
& Y
}
$$
を考える。これは $S$ 上の代数空間の図式である。
$j$ が開埋め込みならば、層の標準的な単射
$j : (X'/Y)_{fin} \to (X/Y)_{fin}$ が存在する。
\end{lemma}

\begin{proof}
$(a, Z)$ が $T$ 上の $X'/Y$ に対する対ならば、
$(a, j(Z))$ は $T$ 上の $X/Y$ に対する対である。
\end{proof}

\begin{lemma}
\label{spaces-more-groupoids-lemma-finite-lives-on-locally-quasi-finite-part}
$S$ をスキームとする。
$f : X \to Y$ を $S$ 上の局所有限型な代数空間の射とする。
$X' \subset X$ を、その上で $f$ が局所準有限となる最大の開部分空間とする。
Morphisms of Spaces, Lemma
\ref{spaces-morphisms-lemma-locally-finite-type-quasi-finite-part} を参照せよ。
このとき $(X/Y)_{fin} = (X'/Y)_{fin}$ である。
\end{lemma}

\begin{proof}
Lemma \ref{spaces-more-groupoids-lemma-finite-open} は単射
$(X'/Y)_{fin} \to (X/Y)_{fin}$ を与える。
Morphisms of Spaces, Lemma
\ref{spaces-morphisms-lemma-locally-finite-type-quasi-finite-part}
により、$X'$ の形成は基底変換と可換である。
したがって、$Z \subset X$ が開部分空間であり、
$f|_Z : Z \to Y$ が有限ならば $Z \subset X'$ であることを証明すればよい。
これは有限射が局所準有限であることから真である。Morphisms of Spaces, Lemma
\ref{spaces-morphisms-lemma-finite-quasi-finite} を参照せよ。
\end{proof}

\begin{lemma}
\label{spaces-more-groupoids-lemma-finite-separated}
$S$ をスキームとする。
$f : X \to Y$ を $S$ 上の代数空間の射とする。
$T$ を $S$ 上の代数空間とし、$(a, Z)$ を
\ref{spaces-more-groupoids-equation-finite-conditions} にある対とする。
$f$ が分離的ならば、$Z$ は $T \times_Y X$ で閉である。
\end{lemma}

\begin{proof}
代数空間の有限射は Morphisms of Spaces, Lemma
\ref{spaces-morphisms-lemma-finite-proper} により普遍的閉である。
$f$ は分離的なので射 $T \times_Y X \to T$ も分離的である。Morphisms of Spaces, Lemma
\ref{spaces-morphisms-lemma-base-change-separated} を参照せよ。
したがって $Z$ の閉性は Morphisms of Spaces, Lemma
\ref{spaces-morphisms-lemma-universally-closed-permanence} から従う。
\end{proof}

\begin{remark}
\label{spaces-more-groupoids-remark-finite-monoid}
$f : X \to Y$ を代数空間の分離的射とする。
層 $(X/Y)_{fin}$ には自然な写像 $(X/Y)_{fin} \to Y$ があり、
これは対 $(a, Z) \in (X/Y)_{fin}(T)$ を元 $a \in Y(T)$ に写す。
Lemma \ref{spaces-more-groupoids-lemma-finite-separated} を用いると、演算
$$
\star_i : (X/Y)_{fin} \times_Y (X/Y)_{fin} \longrightarrow (X/Y)_{fin}
$$
を次の規則で定義できる。
\begin{align*}
\star_1 : ((a, Z_1), (a, Z_2)) & \longmapsto (a, Z_1 \cup Z_2) \\
\star_2 : ((a, Z_1), (a, Z_2)) & \longmapsto (a, Z_1 \cap Z_2) \\
\star_3 : ((a, Z_1), (a, Z_2)) & \longmapsto (a, Z_1 \setminus Z_2) \\
\star_4 : ((a, Z_1), (a, Z_2)) & \longmapsto (a, Z_2 \setminus Z_1).
\end{align*}
これが機能する理由は、$Z_1 \cap Z_2$ が $Z_1$ と $Z_2$ の双方で開かつ閉であることにある
（これはまた、$Z_1 \cup Z_2$ が残りの三つの部分の非交和であることも含意する）。
したがって $(X/Y)_{fin}$ を、$\mathbf{F}_2$-代数
（単位元をもたない）であって $Y$ 上にあるものとみなせる。
その積は $ss' = \star_2(s, s')$ で、和は
$$
s + s' = \star_1(\star_3(s, s'), \star_4(s, s'))
$$
で与えられ、これは結局対称差をとることにほかならない。
この代数の層では $0 = (1_Y, \emptyset)$ であり、実際に
$s + s = 0$ が任意の局所切断 $s$ に対して成り立つことに注意せよ。
$f : X \to Y$ が有限ならば、この代数は単位元 $1 = (1_Y, X)$ をもち、
$\star_3(s, s') = s(1 + s')$ および
$\star_4(s, s') = (1 + s)s'$ が成り立つ。
\end{remark}

\begin{remark}
\label{spaces-more-groupoids-remark-finite-quasi-finite-separated-morphism-schemes}
$f : X \to Y$ をスキームの分離的かつ局所準有限な射とする。
この場合、層 $(X/Y)_{fin}$ は層 $f_!\mathbf{F}_2$ と密接に関係する
% P09-ERR-1367
（将来の参照をここに挿入する）。後者は $Y_\etale$ 上の層である。
すなわち、$V \to Y$ がエタールで $s \in \Gamma(V, f_!\mathbf{F}_2)$ ならば、
$s \in \Gamma(V \times_Y X, \mathbf{F}_2)$ は台
$Z = \text{Supp}(s)$ をもつ切断であり、その台は $V$ 上固有である。
$f$ はさらに局所準有限なので、射影 $Z \to V$ は実際に有限である。
定数アーベル層の切断の台は開であるから、対
$(V \to Y, \text{Supp}(s))$ は \ref{spaces-more-groupoids-equation-finite-conditions} を満たす。
実際、この場合
$f_!\mathbf{F}_2 \cong (X/Y)_{fin}|_{Y_\etale}$
であり、これは $\mathbf{F}_2$-代数構造を説明する。この構造は Remark
\ref{spaces-more-groupoids-remark-finite-monoid} で導入したものである。
\end{remark}

\begin{lemma}
\label{spaces-more-groupoids-lemma-finite-diagonal}
$S$ をスキームとする。
$f : X \to Y$ を $S$ 上の代数空間の射とする。
$(X/Y)_{fin} \to Y$ の対角射
$$
(X/Y)_{fin} \longrightarrow (X/Y)_{fin} \times_Y (X/Y)_{fin}
$$
は（スキームにより）表現可能な開埋め込みであり、「絶対」対角射
$$
(X/Y)_{fin} \longrightarrow (X/Y)_{fin} \times (X/Y)_{fin}
$$
は（スキームにより）表現可能である。
\end{lemma}

\begin{proof}
絶対対角射は、相対対角射と $Y$ の対角射の基底変換との合成であり、
後者はスキームにより表現可能なので、第二の主張は第一の主張から従う。
Spaces, Section \ref{spaces-section-representable} を参照せよ。
第一の主張を証明するには、次を示さなければならない。
スキーム $T$ と二つの対 $(a, Z_1)$、$(a, Z_2)$ が与えられたとする。
これらは $T$ 上にあり、第一成分が同一で \ref{spaces-more-groupoids-equation-finite-conditions} を満たす。
次の性質をもつ開部分スキーム $V \subset T$ が存在する。
スキームの任意の射 $h : T' \to T$ に対して
$$
h(T') \subset V \Leftrightarrow
\Big(T' \times_T Z_1 = T' \times_T Z_2
\text{ as subspaces of }T' \times_Y X\Big)
$$
が成り立つ。
$V$ を構成しよう。$Z_1 \cap Z_2$ は $Z_1$ と $Z_2$ で開であることに注意せよ。
$\text{pr}_0|_{Z_i} : Z_i \to T$ は有限、したがって固有なので
（Morphisms of Spaces, Lemma \ref{spaces-morphisms-lemma-finite-proper} を参照）、
% P09-ERR-1368
$$
E =
\text{pr}_0|_{Z_1}\left(Z_1 \setminus Z_1 \cap Z_2)\right)
\cup
\text{pr}_0|_{Z_2}\left(Z_2 \setminus Z_1 \cap Z_2)\right)
$$
は $T$ で閉である。ここで $V = T \setminus E$ が機能することは明らかである。
\end{proof}

\begin{lemma}
\label{spaces-more-groupoids-lemma-finite-criterion-etale}
$S$ をスキームとする。
$f : X \to Y$ を $S$ 上の代数空間の射とする。
$U$ がスキームで、$U \to Y$ がエタール射であり、
$Z \subset U \times_Y X$ が $U$ 上有限な開部分空間であると仮定する。
このとき誘導射 $U \to (X/Y)_{fin}$ はエタールである。
\end{lemma}

\begin{proof}
これは Lemma \ref{spaces-more-groupoids-lemma-finite-diagonal} にある対角射の記述から形式的に従うが、
理論の展開における重要な段階なので詳述する。
任意のスキーム $T$ をとり、これが $S$ 上にあるとする。
射 $T \to (X/Y)_{fin}$ に対して、射影
$$
T \times_{(X/Y)_{fin}} U \longrightarrow T
$$
がエタールであることを確認しなければならない。次に注意せよ。
$$
T \times_{(X/Y)_{fin}} U
=
(X/Y)_{fin} \times_{((X/Y)_{fin} \times_Y (X/Y)_{fin})} (T \times_Y U)
$$
Lemma \ref{spaces-more-groupoids-lemma-finite-diagonal} の結果を適用すると、
$T \times_{(X/Y)_{fin}} U$ は $T \times_Y U$ の開部分スキームにより表現される。
射影 $T \times_Y U \to T$ は Morphisms of Spaces, Lemma
\ref{spaces-morphisms-lemma-base-change-etale} によりエタールなので、結論を得る。
\end{proof}

\begin{lemma}
\label{spaces-more-groupoids-lemma-finite-pullback}
$S$ をスキームとする。
代数空間の図式
$$
\xymatrix{
X' \ar[d] \ar[r] & X \ar[d] \\
Y' \ar[r] & Y
}
$$
が $S$ 上のファイバー積正方形であるとする。このとき
$$
\xymatrix{
(X'/Y')_{fin} \ar[d] \ar[r] & (X/Y)_{fin} \ar[d] \\
Y' \ar[r] & Y
}
$$
は $(\Sch/S)_{fppf}$ 上の層のファイバー積正方形である。
\end{lemma}

\begin{proof}
定義から直ちに、層 $(X'/Y')_{fin}$ は層
$Y' \times_Y (X/Y)_{fin}$ に等しい。
\end{proof}

\begin{lemma}
\label{spaces-more-groupoids-lemma-finite-surjective-etale-cover}
$S$ をスキームとする。
$f : X \to Y$ を $S$ 上の代数空間の射とする。
$f$ が分離的かつ局所準有限ならば、スキーム $U$ が存在し、これは $Y$ 上エタールであり、
全射エタール射 $U \to (X/Y)_{fin}$ が $Y$ 上に存在する。
\end{lemma}

\begin{proof}
この主張が意味をなすことは、Lemma \ref{spaces-more-groupoids-lemma-finite-diagonal}
にある $(X/Y)_{fin}$ の対角射についての結果から分かる。
Spaces, Lemma \ref{spaces-lemma-representable-diagonal} を参照せよ。
$V$ をスキームとし、$V \to Y$ を全射エタール射とする。
Lemma \ref{spaces-more-groupoids-lemma-finite-pullback} により、射
$(V \times_Y X/V)_{fin} \to (X/Y)_{fin}$ は写像 $V \to Y$ の基底変換であり、
したがって全射かつエタールである。Spaces, Lemma
\ref{spaces-lemma-base-change-representable-transformations-property} を参照せよ。
ゆえに $(V \times_Y X/V)_{fin}$ について補題を証明すれば十分である。
（ここでは、表現可能、全射、かつエタールな関手の変換の合成も、
表現可能、全射、かつエタールであることを暗黙に用いている。Spaces, Lemmas
\ref{spaces-lemma-composition-representable-transformations}、
\ref{spaces-lemma-composition-representable-transformations-property}、および
Morphisms, Lemmas \ref{morphisms-lemma-composition-surjective}、
\ref{morphisms-lemma-composition-etale} を参照せよ。）
分離的および局所準有限という性質は基底変換で保たれることに注意せよ。
Morphisms of Spaces, Lemmas
\ref{spaces-morphisms-lemma-base-change-separated} および
\ref{spaces-morphisms-lemma-base-change-quasi-finite} を参照せよ。
したがって $V \times_Y X \to V$ も分離的かつ局所準有限であり、
Morphisms of Spaces, Proposition
\ref{spaces-morphisms-proposition-locally-quasi-finite-separated-over-scheme}
により $V \times_Y X$ もスキームである。
よって $f : X \to Y$ はスキームの分離的かつ局所準有限な射であると仮定してよい。

\medskip\noindent
点 $y \in Y$ を選び、点 $x_1, \ldots, x_n \in X$ を選ぶ。これらは $y$ 上にある。
エタール近傍 $a : (U, u) \to (Y, y)$ と、More on Morphisms, Lemma
\ref{more-morphisms-lemma-etale-splits-off-quasi-finite-part-technical-variant}
にある分解
% P09-ERR-1369
% P09-ERR-1370
$$
U \times_S X =
W \amalg
\ \coprod\nolimits_{i = 1, \ldots, n}
\ \coprod\nolimits_{j = 1, \ldots, m_j}
V_{i, j}
$$
を選ぶ。任意の部分集合
$$
I \subset \{(i, j) \mid 1 \leq i \leq n, \ 1 \leq j \leq m_i\}.
$$
を選ぶ。これらの選択から、
対 $(a, Z)$ を得る。その第二成分は
$Z = \bigcup_{(i, j) \in I} V_{i, j}$ である。
これは条件 \ref{spaces-more-groupoids-equation-finite-conditions} を満たす。言い換えれば、射
$U \to (X/Y)_{fin}$ を得る。この射の構成は上で選んだすべてのものに依存するので、
本来は
$$
U(y, n, x_1, \ldots, x_n, a, I) \longrightarrow (X/Y)_{fin}
$$
と書くべきである。この射は Lemma \ref{spaces-more-groupoids-lemma-finite-criterion-etale} によりエタールである。

\medskip\noindent
主張：これらすべての非交和は $(X/Y)_{fin}$ 上全射である。
この主張が成り立てば補題が真であることは明らかである。

\medskip\noindent
全射性を示すには次を示さなければならない
（Spaces, Remark \ref{spaces-remark-warning} を参照）。
スキーム $T$ をとり、これが $S$ 上にあるとする。点 $t \in T$、および写像
$T \to (X/Y)_{fin}$ が与えられたとする。上のようなデータ
$(y, n, x_1, \ldots, x_n, a, I)$ であって、$t$ が射影
$$
U(y, n, x_1, \ldots, x_n, a, I) \times_{(X/Y)_{fin}} T \longrightarrow T.
$$
の像に属するものを見いださなければならない。
これを証明するには、$T$ を $\Spec(\overline{\kappa(t)})$ で置き換え、
$T \to (X/Y)_{fin}$ を合成
$\Spec(\overline{\kappa(t)}) \to T \to (X/Y)_{fin}$ で置き換えてよい。
言い換えれば、$T$ は代数閉体のスペクトルであると仮定してよい。

\medskip\noindent
$T = \Spec(k)$ を代数閉体 $k$ のスペクトルとする。
射 $T \to (X/Y)_{fin}$ は、条件 \ref{spaces-more-groupoids-equation-finite-conditions} を満たす対
$(T \to Y, Z)$ により与えられる。図示すると次のようになる。
$$
\xymatrix{
& Z \ar[d] \ar[r] & X \ar[d] \\
\Spec(k) \ar@{=}[r] & T \ar[r] & Y
}
$$
$y \in Y$ を $T \to Y$ の像点とする。
$Z$ は $k$ 上有限なので有限個の点をもつ。
したがって有限個の点 $x_1, \ldots, x_n \in X$ が存在し、
$Z$ の像は、$X$ において $\{x_1, \ldots, x_n\}$ に含まれる。
$a : (U, u) \to (Y, y)$ を選び、これが上のように
$y$ と $x_1, \ldots, x_n$ に適合するようにする。このとき図式
% P09-ERR-1370
$$
\xymatrix{
W \amalg
\ \coprod\nolimits_{i = 1, \ldots, n}
\ \coprod\nolimits_{j = 1, \ldots, m_j}
V_{i, j} \ar[d] \ar[r] &
X \ar[d] \\
U \ar[r] & Y.
}
$$
を得る。$k$ は代数閉で、$\kappa(y) \subset \kappa(u)$ は有限分離なので、
射 $T = \Spec(k) \to Y$ を射
$u = \Spec(\kappa(u)) \to \Spec(\kappa(y)) = y \subset Y$
を通じて因子分解できる。この選択により可換図式
% P09-ERR-1370
$$
\xymatrix{
Z \ar[d] \ar[r] &
W \amalg
\ \coprod\nolimits_{i = 1, \ldots, n}
\ \coprod\nolimits_{j = 1, \ldots, m_j}
V_{i, j} \ar[d] \ar[r] &
X \ar[d] \\
\Spec(k) \ar[r] &
U \ar[r] & Y
}
$$
を得る。左上の矢印の像は $\coprod V_{i, j}$ に入ることが分かっている。
また、定義により $Z$ は $\Spec(k) \times_Y X$ の開部分スキームであり、
これは $(X/Y)_{fin}$ の定義によること、さらに右側の正方形が
ファイバー積正方形であることを想起せよ。したがって
$$
Z \subset
\coprod\nolimits_{i = 1, \ldots, n}\ \coprod\nolimits_{j = 1, \ldots, m_j}
\Spec(k) \times_U V_{i, j}
$$
は開部分スキームである。構成により（More on Morphisms, Lemma
\ref{more-morphisms-lemma-etale-splits-off-quasi-finite-part-technical-variant} を参照）、
各 $V_{i, j}$ は一意な点 $v_{i, j}$ をもち、この点は $u$ 上にあり、剰余体拡大
$\kappa(v_{i, j})/\kappa(u)$ は純非分離である。
したがって各スキーム $\Spec(k) \times_U V_{i, j}$ はちょうど一つの点をもつ。
ゆえに
$$
Z = \coprod\nolimits_{(i, j) \in I} \Spec(k) \times_U V_{i, j}
$$
となる一意な部分集合
$I \subset \{(i, j) \mid 1 \leq i \leq n, \ 1 \leq j \leq m_i\}$
が存在する。定義を展開すると、
$$
U(y, n, x_1, \ldots, x_n, a, I) \times_{(X/Y)_{fin}} T
$$
が、上で得た $I$ に対して空でないことが分かり、所望の結論を得る。
\end{proof}

\begin{proposition}
\label{spaces-more-groupoids-proposition-finite-algebraic-space}
$S$ をスキームとする。
$f : X \to Y$ を $S$ 上の分離的かつ局所有限型な代数空間の射とする。
このとき $(X/Y)_{fin}$ は代数空間である。さらに、射
$(X/Y)_{fin} \to Y$ はエタールである。
\end{proposition}

\begin{proof}
Lemma \ref{spaces-more-groupoids-lemma-finite-lives-on-locally-quasi-finite-part} により、
% P09-ERR-1371
$X$ を $Y$ 上局所準有限な開部分スキームで置き換えてよい。
したがって $f$ は分離的かつ局所準有限であると仮定してよい。
Spaces, Definition \ref{spaces-definition-algebraic-space} の三条件を確認する。
条件 (1) は Lemma \ref{spaces-more-groupoids-lemma-finite-sheaf} から従う。
条件 (2) は Lemma \ref{spaces-more-groupoids-lemma-finite-diagonal} から従う。
最後に条件 (3) は Lemma \ref{spaces-more-groupoids-lemma-finite-surjective-etale-cover} から従う。
したがって $(X/Y)_{fin}$ は代数空間である。
さらに、その補題により可換図式
$$
\xymatrix{
U \ar[rr] \ar[rd] & & (X/Y)_{fin} \ar[ld] \\
& Y
}
$$
が存在する。ここで水平矢印は全射かつエタールであり、南東向きの矢印はエタールである。
Properties of Spaces, Lemma \ref{spaces-properties-lemma-etale-local}
により、南西向きの矢印もエタールである。
\end{proof}

\begin{remark}
\label{spaces-more-groupoids-remark-warning}
$f$ が分離的であるという条件は Proposition
\ref{spaces-more-groupoids-proposition-finite-algebraic-space} から除けない。
例として、$X$ を原点を二重化したアフィン直線とする。Schemes, Example
\ref{schemes-example-affine-space-zero-doubled} を参照せよ。
$Y = \mathbf{A}^1_k$ をアフィン直線とし、$X \to Y$ を明らかな写像とする。
$0 \in Y$ の上には $0_1$ と $0_2$ という二点が $X$ にあることを想起せよ。
したがって $(X/Y)_{fin}$ は $0$ 上に四点、すなわち
$\emptyset, \{0_1\}, \{0_2\}, \{0_1, 0_2\}$ をもつ。
この四点のうち、$U \times_Y X$ の開部分スキームで $U$ 上有限なものへ持ち上げられるのは、
$U \to Y$ がエタールである場合でも三点、すなわち
$\emptyset, \{0_1\}, \{0_2\}$ だけである。
これは $(X/Y)_{fin}$ が代数空間で表現可能だとしても $Y$ 上エタールでないことを示す。
同様の議論により $(X/Y)_{fin}$ は実際には代数空間でないことも分かる。
詳細は省略する。
\end{remark}

\begin{remark}
\label{spaces-more-groupoids-remark-not-scheme}
$Y = \mathbf{A}^1_{\mathbf{R}}$ を実数体上のアフィン直線とし、
$X = \Spec(\mathbf{C})$ を考え、$\mathbf{R}$-有理点 $0$ に写す。この点は $Y$ に属する。
この場合、射 $f : X \to Y$ は有限であるが、$(X/Y)_{fin}$ はスキームではない。
実際、この場合の代数空間 $(X/Y)_{fin}$ は、Spaces, Example
\ref{spaces-example-non-representable-descent} において拡大
$\mathbf{R} \subset \mathbf{C}$ に付随する代数空間と同型であることを示せる。
したがって層 $(X/Y)_{fin}$ を表現するためには、$f$ が有限射であっても、
スキームの圏を離れることが本当に必要である。
\end{remark}

\begin{lemma}
\label{spaces-more-groupoids-lemma-finite-separated-flat-locally-finite-presentation}
$S$ をスキームとする。
$f : X \to Y$ を $S$ 上の分離的、平坦、かつ局所有限表示な代数空間の射とする。
この場合、次が成り立つ。
\begin{enumerate}
\item $(X/Y)_{fin} \to Y$ は分離的、表現可能、かつエタールである。
\item $Y$ がスキームならば、$(X/Y)_{fin}$ はスキーム（により表現可能）である。
\end{enumerate}
\end{lemma}

\begin{proof}
$f$ は特に分離的かつ局所有限型なので（Morphisms of Spaces, Lemma
\ref{spaces-morphisms-lemma-finite-presentation-finite-type} を参照）、
Proposition \ref{spaces-more-groupoids-proposition-finite-algebraic-space} により
$(X/Y)_{fin}$ は代数空間である。
$(X/Y)_{fin} \to Y$ が分離的であることを証明するには、次を示さなければならない。
スキーム $T$ と二つの対 $(a, Z_1)$、$(a, Z_2)$ が与えられ、
これらが $T$ 上にあり、第一成分が同一で条件
\ref{spaces-more-groupoids-equation-finite-conditions} を満たすとする。
次の性質をもつ閉部分スキーム $V \subset T$ が存在する。
スキームの任意の射 $h : T' \to T$ に対して
$$
h \text{ factors through } V \Leftrightarrow
\Big(T' \times_T Z_1 = T' \times_T Z_2
\text{ as subspaces of }T' \times_Y X\Big)
$$
が成り立つ。Lemma \ref{spaces-more-groupoids-lemma-finite-diagonal} の証明では、
% P09-ERR-1372
$V = T' \setminus E$ が $T'$ の開部分スキームであり、その閉補集合が
% P09-ERR-1373
% P09-ERR-1368
$$
E =
\text{pr}_0|_{Z_1}\left(Z_1 \setminus Z_1 \cap Z_2)\right)
\cup
\text{pr}_0|_{Z_2}\left(Z_2 \setminus Z_1 \cap Z_2)\right).
$$
であることを見た。したがって $E$ も開であることを示せばよい。
Lemma \ref{spaces-more-groupoids-lemma-finite-separated} により、$Z_1$ と $Z_2$ は
$T' \times_Y X$ で閉である。ゆえに $Z_1 \setminus Z_1 \cap Z_2$ は $Z_1$ で開である。
$f$ は平坦かつ局所有限表示なので、$\text{pr}_0|_{Z_1}$ もそうである。
これは $Z_1$ が基底変換 $T' \times_Y X$ の開部分空間であることと、
Morphisms of Spaces, Lemmas
\ref{spaces-morphisms-lemma-base-change-finite-presentation} および
% P09-ERR-1374
Lemmas \ref{spaces-morphisms-lemma-base-change-flat} による。
したがって $\text{pr}_0|_{Z_1}$ は開である。Morphisms of Spaces, Lemma
\ref{spaces-morphisms-lemma-fppf-open} を参照せよ。
ゆえに $\text{pr}_0|_{Z_1}\left(Z_1 \setminus Z_1 \cap Z_2)\right)$ は開であり、
$E$ も所望のとおり開である。

\medskip\noindent
$(X/Y)_{fin} \to Y$ がエタールであることはすでに見た。Proposition
\ref{spaces-more-groupoids-proposition-finite-algebraic-space} を参照せよ。
したがって今やこれは局所準有限であり（Morphisms of Spaces, Lemma
\ref{spaces-morphisms-lemma-etale-locally-quasi-finite} を参照）、かつ分離的なので、
Morphisms of Spaces, Lemma
\ref{spaces-morphisms-lemma-locally-quasi-finite-separated-representable}
により表現可能である。最後の主張は明らかである
（必要なら Morphisms of Spaces, Proposition
\ref{spaces-morphisms-proposition-locally-quasi-finite-separated-over-scheme} を用いよ）。
\end{proof}

\noindent
変種：$S$ をスキームとする。
$f : X \to Y$ を $S$ 上の代数空間の射とする。
$\sigma : Y \to X$ を $f$ の切断とする。
代数空間またはスキーム $T$ をとり、これは $S$ 上にあるとする。
次を満たす対 $(a, Z)$ を考える。
\begin{equation}
\label{spaces-more-groupoids-equation-finite-conditions-variant}
\begin{matrix}
a : T \to Y\text{ is a morphism over }S, \\
Z \subset T \times_Y X\text{ is an open subspace} \\
\text{such that }\text{pr}_0|_Z : Z \to T\text{ is finite and} \\
(1_T, \sigma \circ a) : T \to T \times_Y X\text{ factors through }Z.
\end{matrix}
\end{equation}
% P09-ERR-1375
$(X/Y, \sigma)_{fin}$ を、これらの対をパラメータ付ける
$(X/Y)_{fin}$ の部分関手と記す。

\begin{lemma}
\label{spaces-more-groupoids-lemma-finite-plus-section}
$S$ をスキームとする。
$f : X \to Y$ を $S$ 上の代数空間の射とする。
$\sigma : Y \to X$ を $f$ の切断とする。上で定義した関手の変換
$$
t : (X/Y, \sigma)_{fin} \longrightarrow (X/Y)_{fin}.
$$
を考える。このとき次が成り立つ。
\begin{enumerate}
\item $t$ は開埋め込みにより表現可能である。
\item $f$ が分離的ならば、$t$ は開かつ閉な埋め込みにより表現可能である。
\item $(X/Y)_{fin}$ が代数空間ならば、$(X/Y, \sigma)_{fin}$ は代数空間であり、
$(X/Y)_{fin}$ の開部分空間である。
\item $(X/Y)_{fin}$ がスキームならば、$(X/Y, \sigma)_{fin}$ はその開部分スキームである。
\end{enumerate}
\end{lemma}

\begin{proof}
省略する。ヒント：$(a, Z)$ を (\ref{spaces-more-groupoids-equation-finite-conditions}) にある $T$ 上の対とすると、
$Z$ の逆像を考える。この逆像は
$(1_T, \sigma \circ a) : T \to T \times_Y X$ によるものであり、
求める $T$ の開部分スキームである。
\end{proof}





\section{射の有限集合}
\label{spaces-more-groupoids-section-finite-set-arrows}

\noindent
$\mathcal{C}$ を亜群とする。Categories, Definition
\ref{categories-definition-groupoid} を参照せよ。
Groupoids, Section \ref{groupoids-section-groupoids}
で論じたように、これは 7 つ組
$(\text{Ob}, \text{Arrows}, s, t, c, e, i)$ に対応する。

\medskip\noindent
このデータを用いて、別の亜群 $\mathcal{C}_{fin}$ を
次のように作ることができる。
\begin{enumerate}
\item $\mathcal{C}_{fin}$ の対象は、次の性質をもつ有限部分集合
$Z \subset \text{Arrows}$ からなる。
\begin{enumerate}
\item $s(Z) = \{u\}$ は一点集合であり、
\item $e(u) \in Z$ である。
\end{enumerate}
\item $\mathcal{C}_{fin}$ の射は対
$(Z, z)$ からなり、ここで $Z$ は $\mathcal{C}_{fin}$ の対象であり、
$z \in Z$ である。
\item $(Z, z)$ の始域は $Z$ である。
\item $(Z, z)$ の終域は
$t(Z, z) = \{z' \circ z^{-1}; z' \in Z\}$ である。
\item $(Z_1, z_1)$、$(Z_2, z_2)$ が与えられ、
$s(Z_1, z_1) = t(Z_2, z_2)$ を満たすとき、合成
$(Z_1, z_1) \circ (Z_2, z_2)$ は $(Z_2, z_1 \circ z_2)$ である。
\end{enumerate}
これが亜群を定めることの確認は省略する。
図式的には、$\mathcal{C}_{fin}$ の対象は
次の図として見ることができる。
$$
\xymatrix{
& \bullet \\
\bullet \ar@(ul, dl)[]_e \ar[ru] \ar[r] \ar[rd] & \bullet \\
& \bullet
}
$$
$\mathcal{C}_{fin}$ の射を作るには、射の一つを選び、
その逆射を他の射の前に合成する。例えば中央の水平な射を選ぶと、
終域は次の図になる。
$$
\xymatrix{
& \bullet \\
\bullet & \bullet \ar[l] \ar[u] \ar@(dr, ur)[]_e \ar[d] \\
& \bullet
}
$$
$s(Z, z)$ と $t(Z, z)$ の濃度は等しいことに注意せよ。
したがって $\mathcal{C}_{fin}$ は、実際には亜群の可算な非交和である。




\section{亜群の有限部分}
\label{spaces-more-groupoids-section-finite-part-groupoid}

\noindent
この節では、Section \ref{spaces-more-groupoids-section-finite-set-arrows}
で説明した着想を用いて、代数空間の亜群の有限部分を取り出す。

\medskip\noindent
$S$ をスキームとする。
$B$ を $S$ 上の代数空間とする。
$(U, R, s, t, c, e, i)$ を $B$ 上の代数空間の亜群とする。
仮定：射 $s, t$ は分離的かつ局所有限型である。
% P09-ERR-1376
この記法と仮定を、この節を通じて固定する。

\medskip\noindent
$R_s$ により、代数空間 $R$ を $U$ 上の代数空間とみなし、
その構造射を $s$ としたものを表す。$U' = (R_s/U, e)_{fin}$ とおく。
$s$ は分離的かつ局所有限型なので、
Proposition \ref{spaces-more-groupoids-proposition-finite-algebraic-space} および
Lemma \ref{spaces-more-groupoids-lemma-finite-plus-section} により、
$U'$ はエタール射 $g : U' \to U$ を備えた代数空間である。
さらに、Lemma \ref{spaces-more-groupoids-lemma-finite-sheaf} により、
普遍的開部分空間
$Z_{univ} \subset R \times_{s, U, g} U'$ が存在し、これは $U'$ 上有限で、
$(1_{U'}, e \circ g) : U' \to R \times_{s, U, g} U'$
は $Z_{univ}$ を経由する。さらに、
Lemma \ref{spaces-more-groupoids-lemma-finite-separated} により、
% P09-ERR-1377
開部分空間 $Z_{univ}$ は $R \times_{s, U', g} U$ の中で閉でもある。
ここまでの図は次のとおりである。
$$
\xymatrix{
Z_{univ} \ar[d] \ar[rd] & \\
R \times_{s, U, g} U' \ar[d] \ar[r] & U' \ar[d]^g \\
R \ar[r]^s & U
}
$$
$T$ を $B$ 上のスキームとする。$T$ 値点としての $Z_{univ}$ の点は、
次を満たす 3 つ組 $(u, Z, z)$ とみなせる。
\begin{enumerate}
\item $u : T \to U$ は $T$ 値点としての $U$ の点であり、
\item $Z \subset R \times_{s, U, u} T$ は $T$ 上有限な開かつ閉な部分空間で、
$(e \circ u, 1_T)$ がこれを経由し、
\item $z : T \to R$ は $T$ 値点としての $R$ の点で、$s \circ z = u$ を満たし、
$(z, 1_T)$ が $Z$ を経由する。
\end{enumerate}
以上を踏まえると、Section \ref{spaces-more-groupoids-section-finite-set-arrows}
の議論から、$(Z_{univ}, U')$ を $B$ 上の代数空間の亜群にできることは
直観的に明らかである。
% P09-ERR-1378
確実を期すため、関手的な見方を用いて射 $s', t', c', e', i'$ を
一つずつ定義する。（Section \ref{spaces-more-groupoids-section-finite-set-arrows}
の簡単な構成を読み、理解する前に以下を読まないでほしい。）

\medskip\noindent
射 $s' : Z_{univ} \to U'$ は次の規則に対応する。
$$
s' : (u, Z, z) \mapsto (u, Z).
$$
射 $t' : Z_{univ} \to U'$ は次の規則で与えられる。
$$
t' : (u, Z, z) \mapsto (t \circ z, c(Z, i \circ z)).
$$
項 $c(Z, i \circ z)$ は意味をもつ。実際、写像
$c(-, i \circ z) : R \times_{s, U, u} T \to R \times_{s, U, t \circ z} T$
は $c(-, z)$ を逆写像とする同型である。
射 $e' : U' \to Z_{univ}$ は次の規則で与えられる。
$$
e' : (u, Z) \mapsto (u, Z, (e \circ u, 1_T)).
$$
これは、$(e \circ u, 1_T)$ が $Z$ を経由するという要請により
意味をもつことに注意せよ。
射 $i' : Z_{univ} \to Z_{univ}$ は次の規則で与えられる。
$$
i' : (u, Z, z) \mapsto (t \circ z, c(Z, i \circ z), i \circ z).
$$
最後に、合成を次の規則で定義する。
$$
c' : ((u_1, Z_1, z_1), (u_2, Z_2, z_2)) \mapsto (u_2, Z_2, z_1 \circ z_2).
$$
$(U', Z_{univ}, s', t', c', e', i')$ について
代数空間の亜群の公理が成り立つことの確認は省略する。

\medskip\noindent
最後の情報として、亜群の標準射
$$
(U', Z_{univ}, s', t', c', e', i')
\longrightarrow
(U, R, s, t, c, e, i)
$$
がある。実際、射 $U' \to U$ は $g : U' \to U$ であり、
規則 $(u, Z) \mapsto u$ によって定義される。
射 $Z_{univ} \to R$ は規則 $(u, Z, z) \mapsto z$ によって定義される。
これで構成は完了した。得られたことを次のようにまとめる。

\begin{lemma}
\label{spaces-more-groupoids-lemma-finite-part-groupoid}
$S$ をスキームとする。
$B$ を $S$ 上の代数空間とする。
$(U, R, s, t, c, e, i)$ を $B$ 上の代数空間の亜群とする。
射 $s, t$ は分離的かつ局所有限型であると仮定する。
標準射
$$
(U', Z_{univ}, s', t', c', e', i')
\longrightarrow
(U, R, s, t, c, e, i)
$$
が代数空間の亜群の間に存在し、これは $B$ 上の射である。ここで
\begin{enumerate}
\item $g : U' \to U$ は $(R_s/U, e)_{fin} \to U$ と同一視され、
\item $Z_{univ} \subset R \times_{s, U, g} U'$ は普遍的な開（かつ閉）部分空間で、
$U'$ 上有限であり、単位射 $e$ の基底変換を含む。
\end{enumerate}
\end{lemma}

\begin{proof}
上の議論を参照せよ。
\end{proof}









\section{亜群スキームのエタール局所化}
\label{spaces-more-groupoids-section-etale-localize}

\noindent
この節では、\cite[Proposition 4.2]{K-M} と同様の結果を証明する。
ここではもう少し一般的にし、Hilbert スキームを使う代わりに
射の有限部分を用いることで、その使用を避けることを試みる。
目標は、エタール局所化の後で、一点上の代数空間の亜群を
「分裂」させることである。定義は次のとおりである
（\cite[Definition 4.1]{K-M} と非常によく似ている）。

\begin{definition}
\label{spaces-more-groupoids-definition-split-at-point}
$S$ をスキームとする。$B$ を $S$ 上の代数空間とする。
% P09-ERR-1379
$(U, R, s, t, c)$ を $B$ 上の代数空間の亜群とする。
$u \in |U|$ を一点とする。
\begin{enumerate}
\item $R$ が $u$ 上で {\it 強分裂する} というのは、
次を満たす開部分空間 $P \subset R$ が存在するときである。
\begin{enumerate}
\item $(U, P, s|_P, t|_P, c|_{P \times_{s, U, t} P})$ は
$B$ 上の代数空間の亜群であり、
\item $s|_P$, $t|_P$ は有限であり、
\item $\{r \in |R| : s(r) = u, t(r) = u\} \subset |P|$ である。
\end{enumerate}
このような $P$ の選択を、$R$ の $u$ 上の
{\it 強分裂} と呼ぶ。
\item $R$ が $u$ 上で {\it 分裂する} というのは、
次を満たす開部分空間 $P \subset R$ が存在するときである。
\begin{enumerate}
\item $(U, P, s|_P, t|_P, c|_{P \times_{s, U, t} P})$ は
$B$ 上の代数空間の亜群であり、
\item $s|_P$, $t|_P$ は有限であり、
\item $\{g \in |G| : g\text{ maps to }u\} \subset |P|$ である。
ここで $G \to U$ は固定化群である。
\end{enumerate}
このような $P$ の選択を、$R$ の $u$ 上の
{\it 分裂} と呼ぶ。
\item $R$ が $u$ 上で {\it 準分裂する} というのは、
次を満たす開部分空間 $P \subset R$ が存在するときである。
\begin{enumerate}
\item $(U, P, s|_P, t|_P, c|_{P \times_{s, U, t} P})$ は
$B$ 上の代数空間の亜群であり、
\item $s|_P$, $t|_P$ は有限であり、
\item $e(u) \in |P|$\footnote{この条件は (a) から従う。}。
\end{enumerate}
このような $P$ の選択を、$R$ の $u$ 上の {\it 準分裂} と呼ぶ。
\end{enumerate}
\end{definition}

\noindent
$P$ に課した条件と (\ref{spaces-more-groupoids-equation-finite-conditions}) の対に課した条件が
似ていることに注意せよ。特に $s, t$ が分離的ならば、
$P$ は $R$ の中で閉でもある
（Lemma \ref{spaces-more-groupoids-lemma-finite-separated} を参照）。

\medskip\noindent
代数空間の亜群 $(U, R, s, t, c)$ が $B$ 上にあり、さらに
一点 $u \in |U|$ から出発するとする。
目標はエタール局所化後に亜群を分裂させることなので、
$U$ をアフィンスキームで置き換えてもよい
（これは考えられるすべての応用に無害だという意味である）。
さらに、これから課す追加仮定により、$R$ は少なくとも
$\{r \in |R| : s(r) = u, t(r) = u\}$ または $e(u)$ の近傍で
スキームになる。このため、以下に述べる亜群スキームから始める。
しかし証明の手法によってスキームの圏の外へ出るので、
上では代数空間の亜群の場合について分裂を定式化した。
一方、射 $s$, $t$ がさらに平坦かつ有限表示である場合以外の
応用は知られておらず、その場合は最終的に再びスキームの圏へ戻る。

\begin{situation}[強分裂]
\label{spaces-more-groupoids-situation-strong-splitting}
$S$ をスキームとする。
$(U, R, s, t, c)$ を $S$ 上の亜群スキームとする。
$u \in U$ を一点とする。次を仮定する。
\begin{enumerate}
\item $s, t : R \to U$ は分離的であり、
\item $s$, $t$ は局所有限型であり、
\item 集合 $\{r \in R : s(r) = u, t(r) = u\}$ は有限であり、
\item $s$ は (3) の集合の各点で準有限である。
\end{enumerate}
仮定 (3) と (4) は、ファイバー $s^{-1}(\{u\})$ が有限であるという
仮定から従うことに注意せよ。Morphisms, Lemma
\ref{morphisms-lemma-finite-fibre} を参照せよ。
\end{situation}

\begin{situation}[分裂]
\label{spaces-more-groupoids-situation-splitting}
$S$ をスキームとする。
$(U, R, s, t, c)$ を $S$ 上の亜群スキームとする。
$u \in U$ を一点とする。次を仮定する。
\begin{enumerate}
\item $s, t : R \to U$ は分離的であり、
\item $s$, $t$ は局所有限型であり、
\item 集合 $\{g \in G : g\text{ maps to }u\}$ は有限である。
ここで $G \to U$ は固定化群であり、
\item $s$ は (3) の集合の各点で準有限である。
\end{enumerate}
\end{situation}

\begin{situation}[準分裂]
\label{spaces-more-groupoids-situation-quasi-splitting}
$S$ をスキームとする。
$(U, R, s, t, c)$ を $S$ 上の亜群スキームとする。
$u \in U$ を一点とする。次を仮定する。
\begin{enumerate}
\item $s, t : R \to U$ は分離的であり、
\item $s$, $t$ は局所有限型であり、
\item $s$ は $e(u)$ で準有限である。
\end{enumerate}
\end{situation}

\noindent
代数空間の存在定理への応用には、準分裂の場合で十分である。
さらに、準分裂の場合から準 DM スタックのエタール局所構造定理を
証明できる。分裂の場合は Keel--Mori 定理の一つの版を証明するために
用いる。強分裂の場合からは、準コンパクトな対角射をもつ
準 DM 代数スタックのエタール局所構造定理が得られる。

\begin{lemma}[強分裂の存在]
\label{spaces-more-groupoids-lemma-strong-splitting}
Situation \ref{spaces-more-groupoids-situation-strong-splitting} において、代数空間 $U'$、
エタール射 $U' \to U$、および
点 $u' : \Spec(\kappa(u)) \to U'$ が存在し、これは
$u : \Spec(\kappa(u)) \to U$ の上にある。また、制限
$R' = R|_{U'}$、すなわち $R$ の $U'$ への制限は $u'$ 上で強分裂する。
\end{lemma}

\begin{proof}
Lemma \ref{spaces-more-groupoids-lemma-finite-part-groupoid} で構成した
$f : (U', Z_{univ}, s', t', c') \to (U, R, s, t, c)$ をとる。
$R' = R \times_{(U \times_S U)} (U' \times_S U')$ であることを思い出す。
したがって、代数空間の亜群の射
$(f, t', s') : Z_{univ} \to R'$ が得られる。
$$
(U', Z_{univ}, s', t', c') \to (U', R', s', t', c')
$$
（記法を濫用し、二つの亜群の射を同じ記号で表している。）
さて、$Z_{univ} \subset R \times_{s, U, g} U'$ は開であり、
$R' \to R \times_{s, U, g} U'$ は
$U' \to U$ の基底変換としてエタールなので、
$Z_{univ} \to R'$ は開埋め込みである。
構成により、射 $s', t' : Z_{univ} \to U'$ は有限である。
残るのは点 $u'$ を $U'$ に見つけることである。

\medskip\noindent
補題の主張と同様に、$u$ を射 $\Spec(\kappa(u)) \to U$ とみなす。
$F_u = R \times_{s, U} \Spec(\kappa(u))$ とおく。
集合 $\{r \in R : s(r) = u, t(r) = u\}$ は仮定により有限であり、
$F_u \to \Spec(\kappa(u))$ は仮定によりその各元で準有限である。
したがって、開かつ閉な部分スキームへの分解
$$
F_u = Z_u \amalg Rest
$$
をとることができる。ここでスキーム $Z_u$ は $\kappa(u)$ 上有限で、
その台は $\{r \in R : s(r) = u, t(r) = u\}$ である。
$e(u) \in Z_u$ に注意せよ。したがって
Section \ref{spaces-more-groupoids-section-finite-part-groupoid} における $U'$ の構成により、
$(u, Z_u)$ は $\Spec(\kappa(u))$ 値点 $u'$ を $U'$ 上に定める。

\medskip\noindent
集合
$\{r' \in |R'| : s'(r') = u', t'(r') = u'\}$
が $|Z_{univ}|$ に含まれることを示す必要がまだある。
この集合の任意の点 $r'$ をとり、射 $z' : \Spec(k) \to R'$ で表す。
$z : \Spec(k) \to R$ と記す。これは
$z'$ と写像 $R' \to R$ の合成である。
明らかに $z$ は集合
$\{r \in R : s(r) = u, t(r) = u\}$ の元を定める。
また、合成 $s \circ z, t \circ z : \Spec(k) \to U$ は
$u$ を経由するので、$s \circ z, t \circ z$ を
射 $\Spec(k) \to \Spec(\kappa(u))$ とみなせる。このとき、
% P09-ERR-1380
$z' = (z, u' \circ t \circ z, u'\circ s \circ u)$ であり、
これは $R' = R \times_{(U \times_S U)} (U' \times_S U')$
への射としての等式である。
3 つ組
$$
(s \circ z, Z_u \times_{\Spec(\kappa(u)), s \circ z} \Spec(k), z)
$$
を考える。ここで $Z_u$ は上のものである。これは
$\Spec(k)$ 値点としての $Z_{univ}$ の点を定め、
$s', t'$ による $U'$ への像は $u'$ であり、
$Z_{univ} \to R'$ による像は点 $r'$ である。これは、上で述べた
$z$ と $z'$ の関係による。
これで証明は完了する。
\end{proof}

\begin{lemma}[分裂の存在]
\label{spaces-more-groupoids-lemma-splitting}
Situation \ref{spaces-more-groupoids-situation-splitting} において、代数空間 $U'$、
エタール射 $U' \to U$、および
点 $u' : \Spec(\kappa(u)) \to U'$ が存在し、これは
$u : \Spec(\kappa(u)) \to U$ の上にある。また、制限
$R' = R|_{U'}$、すなわち $R$ の $U'$ への制限は $u'$ 上で分裂する。
\end{lemma}

\begin{proof}
Lemma \ref{spaces-more-groupoids-lemma-finite-part-groupoid} で構成した
$f : (U', Z_{univ}, s', t', c') \to (U, R, s, t, c)$ をとる。
$R' = R \times_{(U \times_S U)} (U' \times_S U')$ であることを思い出す。
したがって、代数空間の亜群の射
$(f, t', s') : Z_{univ} \to R'$ が得られる。
$$
(U', Z_{univ}, s', t', c') \to (U', R', s', t', c')
$$
（記法を濫用し、二つの亜群の射を同じ記号で表している。）
さて、$Z_{univ} \subset R \times_{s, U, g} U'$ は開であり、
$R' \to R \times_{s, U, g} U'$ は
$U' \to U$ の基底変換としてエタールなので、
$Z_{univ} \to R'$ は開埋め込みである。
構成により、射 $s', t' : Z_{univ} \to U'$ は有限である。
残るのは点 $u'$ を $U'$ に見つけることである。

\medskip\noindent
補題の主張と同様に、$u$ を射 $\Spec(\kappa(u)) \to U$ とみなす。
$F_u = R \times_{s, U} \Spec(\kappa(u))$ とおく。
$G_u \subset F_u$ を、$G \to U$ の $u$ 上の
スキーム論的ファイバーとする。仮定により $G_u$ は有限であり、
$F_u \to \Spec(\kappa(u))$ は仮定により $G_u$ の各点で準有限である。
したがって、開かつ閉な部分スキームへの分解
$$
F_u = Z_u \amalg Rest
$$
をとることができる。ここでスキーム $Z_u$ は $\kappa(u)$ 上有限で、
その台は $G_u$ である。$e(u) \in Z_u$ に注意せよ。
したがって Section \ref{spaces-more-groupoids-section-finite-part-groupoid} における
$U'$ の構成により、$(u, Z_u)$ は
$\Spec(\kappa(u))$ 値点 $u'$ を $U'$ 上に定める。

\medskip\noindent
集合 $\{g' \in |G'| : g'\text{ maps to }u'\}$ が
$|Z_{univ}|$ に含まれることを示す必要がまだある。
この集合の任意の点 $g'$ をとり、
射 $z' : \Spec(k) \to G'$ で表す。
$z : \Spec(k) \to G$ と記す。これは
$z'$ と写像 $G' \to G$ の合成である。
明らかに $z$ は $G_u$ の点を定める。実際、対応する
$\tilde u : \Spec(k) \to u \to U$ と書くことにする。これは
$u$ または $U$ への写像である。
3 つ組
$$
(\tilde u, Z_u \times_{u, \tilde u} \Spec(k), z)
$$
を考える。ここで $Z_u$ は上のものである。これは
$\Spec(k)$ 値点としての $Z_{univ}$ の点を定め、
$s', t'$ による $U'$ への像は $u'$ であり、
$Z_{univ} \to R'$ による像は点 $z'$ である
（$R$ における像が $z$ だからである）。
これで証明は完了する。
\end{proof}

\begin{lemma}[準分裂の存在]
\label{spaces-more-groupoids-lemma-quasi-splitting}
Situation \ref{spaces-more-groupoids-situation-quasi-splitting} において、代数空間 $U'$、
エタール射 $U' \to U$、および点
$u' : \Spec(\kappa(u)) \to U'$ が存在し、これは
$u : \Spec(\kappa(u)) \to U$ の上にある。また、制限
$R' = R|_{U'}$、すなわち $R$ の $U'$ への制限は $u'$ 上で準分裂する。
\end{lemma}

\begin{proof}
Lemma \ref{spaces-more-groupoids-lemma-finite-part-groupoid} で構成した
$f : (U', Z_{univ}, s', t', c') \to (U, R, s, t, c)$ をとる。
$R' = R \times_{(U \times_S U)} (U' \times_S U')$ であることを思い出す。
したがって、代数空間の亜群の射
$(f, t', s') : Z_{univ} \to R'$ が得られる。
$$
(U', Z_{univ}, s', t', c') \to (U', R', s', t', c')
$$
（記法を濫用し、二つの亜群の射を同じ記号で表している。）
さて、$Z_{univ} \subset R \times_{s, U, g} U'$ は開であり、
$R' \to R \times_{s, U, g} U'$ は
$U' \to U$ の基底変換としてエタールなので、
$Z_{univ} \to R'$ は開埋め込みである。
構成により、射 $s', t' : Z_{univ} \to U'$ は有限である。
残るのは点 $u'$ を $U'$ に見つけることである。

\medskip\noindent
補題の主張と同様に、$u$ を射 $\Spec(\kappa(u)) \to U$ とみなす。
$F_u = R \times_{s, U} \Spec(\kappa(u))$ とおく。
射 $F_u \to \Spec(\kappa(u))$ は、仮定により $e(u)$ で準有限である。
したがって、開かつ閉な部分スキームへの分解
$$
F_u = Z_u \amalg Rest
$$
をとることができる。ここでスキーム $Z_u$ は $\kappa(u)$ 上有限で、
その台は $e(u)$ である。
したがって Section \ref{spaces-more-groupoids-section-finite-part-groupoid} における
$U'$ の構成により、$(u, Z_u)$ は
$\Spec(\kappa(u))$ 値点 $u'$ を $U'$ 上に定める。
証明を終えるには $e'(u') \in Z_{univ}$ を示さなければならないが、
これは明らかである。
\end{proof}

\noindent
最後に、追加の仮定を置くとスキームが得られる。

\begin{lemma}
\label{spaces-more-groupoids-lemma-strong-splitting-scheme}
Situation \ref{spaces-more-groupoids-situation-strong-splitting} において、さらに
$s, t$ は平坦かつ局所有限表示であると仮定する。
このときスキーム $U'$、分離エタール射
$U' \to U$、および点 $u' \in U'$ が存在し、この点は $u$ の上にあり、
$\kappa(u) = \kappa(u')$ を満たす。また、制限
$R' = R|_{U'}$、すなわち $R$ の $U'$ への制限は
$u'$ 上で強分裂する。
\end{lemma}

\begin{proof}
これは Lemma \ref{spaces-more-groupoids-lemma-strong-splitting} の証明における
$U'$ の構成から従う。実際、この場合
$U' = (R_s/U, e)_{fin}$ はスキームであり、
Lemmas \ref{spaces-more-groupoids-lemma-finite-separated-flat-locally-finite-presentation} および
\ref{spaces-more-groupoids-lemma-finite-plus-section} により $U$ 上分離的である。
\end{proof}

\begin{lemma}
\label{spaces-more-groupoids-lemma-splitting-scheme}
Situation \ref{spaces-more-groupoids-situation-splitting} において、さらに
$s, t$ は平坦かつ局所有限表示であると仮定する。
このときスキーム $U'$、分離エタール射
$U' \to U$、および点 $u' \in U'$ が存在し、この点は $u$ の上にあり、
$\kappa(u) = \kappa(u')$ を満たす。また、制限
$R' = R|_{U'}$、すなわち $R$ の $U'$ への制限は
$u'$ 上で分裂する。
\end{lemma}

\begin{proof}
これは Lemma \ref{spaces-more-groupoids-lemma-splitting} の証明における
$U'$ の構成から従う。実際、この場合
$U' = (R_s/U, e)_{fin}$ はスキームであり、
Lemmas \ref{spaces-more-groupoids-lemma-finite-separated-flat-locally-finite-presentation} および
\ref{spaces-more-groupoids-lemma-finite-plus-section} により $U$ 上分離的である。
\end{proof}

\begin{lemma}
\label{spaces-more-groupoids-lemma-quasi-splitting-scheme}
Situation \ref{spaces-more-groupoids-situation-quasi-splitting} において、さらに
$s, t$ は平坦かつ局所有限表示であると仮定する。
このときスキーム $U'$、分離エタール射
$U' \to U$、および $u' \in U'$ が存在し、これは $u$ の上にあり、
$\kappa(u) = \kappa(u')$ を満たす。また、制限
$R' = R|_{U'}$、すなわち $R$ の $U'$ への制限は
$u'$ 上で準分裂する。
\end{lemma}

\begin{proof}
これは Lemma \ref{spaces-more-groupoids-lemma-quasi-splitting} の証明における
$U'$ の構成から従う。実際、この場合
$U' = (R_s/U, e)_{fin}$ はスキームであり、
Lemmas \ref{spaces-more-groupoids-lemma-finite-separated-flat-locally-finite-presentation} および
\ref{spaces-more-groupoids-lemma-finite-plus-section} により $U$ 上分離的である。
\end{proof}

\noindent
実際、有限局所自由な亜群に関する先の結果を適用すれば、
アフィンスキームを得ることができる。

\begin{lemma}
\label{spaces-more-groupoids-lemma-strong-splitting-affine-scheme}
Situation \ref{spaces-more-groupoids-situation-strong-splitting} において、さらに
$s, t$ は平坦かつ局所有限表示であり、$U$ はアフィンであると仮定する。
このときアフィンスキーム $U'$、エタール射
$U' \to U$、および $u' \in U'$ が存在し、これは $u$ の上にあり、
$\kappa(u) = \kappa(u')$ を満たす。また、制限
$R' = R|_{U'}$、すなわち $R$ の $U'$ への制限は
$u'$ 上で強分裂する。
\end{lemma}

\begin{proof}
Lemma \ref{spaces-more-groupoids-lemma-strong-splitting-scheme} で得た、スキームの分離エタール射
$U' \to U$ と点 $u' \in U'$ をとる。
$P \subset R'$ を、$R'$ の $u'$ 上の強分裂とする。
More on Groupoids, Lemma
\ref{more-groupoids-lemma-restrict-preserves-type} により、射
$s', t' : R' \to U'$ は平坦かつ局所有限表示である。
% P09-ERR-1381
これらは仮定により有限である。したがって $s', t'$ は有限局所自由である。
Morphisms, Lemma \ref{morphisms-lemma-finite-flat} を参照せよ。
% P09-ERR-1382
特に $t(s^{-1}(u'))$ は有限な点集合
$\{u'_1, u'_2, \ldots, u'_n\}$ であり、これらは $U'$ の点である。
準コンパクトな開部分 $W \subset U'$ を選び、各 $u'_i$ を含むものとする。
$U$ はアフィンなので、射
$W \to U$ は準コンパクトである
（Schemes, Lemma \ref{schemes-lemma-quasi-compact-affine} を参照）。
射 $W \to U$ はさらに局所準有限
（Morphisms, Lemma \ref{morphisms-lemma-etale-locally-quasi-finite} を参照）
かつ分離的である。したがって More on Morphisms,
Lemma \ref{more-morphisms-lemma-quasi-finite-separated-quasi-affine}
（Zariski の主定理の一つの版）により、$W$ は準アフィンである。
Properties, Lemma \ref{properties-lemma-ample-finite-set-in-affine} により、
% P09-ERR-1383
$\{u'_1, \ldots, u'_n\}$ は $U'$ のアフィン開部分に含まれる。
したがって Groupoids, Lemma \ref{groupoids-lemma-find-invariant-affine}
を適用すると、アフィンな $P$ 不変開部分
$U'' \subset U'$ が存在し、これは $u'$ を含む。

\medskip\noindent
% P09-ERR-1384
証明を終えるため、$R'' = R|_{U''}$ を $R$ の $U''$ への制限と記す。
これは $R'$ の $U''$ への制限と同じである。
$P \subset R'$ は開かつ閉な部分スキームなので、
$P|_{U''} \subset R''$ もそうである。構成により開部分スキーム
$U'' \subset U'$ は $P$ 不変であり、これは
$$
P|_{U''} = (s'|_P)^{-1}(U'') = (t'|_P)^{-1}(U'')
$$
を意味する（Groupoids, Section
\ref{groupoids-section-invariant} の議論を参照）。
したがって、$s''$ と $t''$ の $P|_{U''}$ への制限は依然として有限である。
% P09-ERR-1385
部分亜群スキーム $P|_{U''}$ は依然として $R''$ の $u''$ 上の強分裂である。
上では (a)、(b)、(c) が成り立つことを確認した。実際、
$$
\{r' \in R': t'(r') = u', s'(r') = u'\} =
\{r'' \in R'': t''(r'') = u', s''(r'') = u'\}
$$
は自明である。
これで補題は証明された。
\end{proof}

\begin{lemma}
\label{spaces-more-groupoids-lemma-splitting-affine-scheme}
Situation \ref{spaces-more-groupoids-situation-splitting} において、さらに
$s, t$ は平坦かつ局所有限表示であり、$U$ はアフィンであると仮定する。
このときアフィンスキーム $U'$、エタール射
$U' \to U$、および $u' \in U'$ が存在し、これは $u$ の上にあり、
$\kappa(u) = \kappa(u')$ を満たす。また、制限
$R' = R|_{U'}$、すなわち $R$ の $U'$ への制限は
$u'$ 上で分裂する。
\end{lemma}

\begin{proof}
この補題の証明は Lemma
\ref{spaces-more-groupoids-lemma-strong-splitting-affine-scheme} の証明と文字どおり同じである。
ただし「強分裂」を「分裂」で置き換え（2 回）、
Lemma \ref{spaces-more-groupoids-lemma-strong-splitting-scheme} への参照を
Lemma \ref{spaces-more-groupoids-lemma-splitting-scheme} への参照で置き換える必要がある。
\end{proof}

\begin{lemma}
\label{spaces-more-groupoids-lemma-quasi-splitting-affine-scheme}
Situation \ref{spaces-more-groupoids-situation-quasi-splitting} において、さらに
$s, t$ は平坦かつ局所有限表示であり、$U$ はアフィンであると仮定する。
このときアフィンスキーム $U'$、エタール射
$U' \to U$、および $u' \in U'$ が存在し、これは $u$ の上にあり、
$\kappa(u) = \kappa(u')$ を満たす。また、制限
$R' = R|_{U'}$、すなわち $R$ の $U'$ への制限は
$u'$ 上で準分裂する。
\end{lemma}

\begin{proof}
この補題の証明は Lemma
\ref{spaces-more-groupoids-lemma-strong-splitting-affine-scheme} の証明と文字どおり同じである。
ただし「強分裂」を「準分裂」で置き換え（2 回）、
Lemma \ref{spaces-more-groupoids-lemma-strong-splitting-scheme} への参照を
Lemma \ref{spaces-more-groupoids-lemma-quasi-splitting-scheme} への参照で置き換える必要がある。
\end{proof}
